\documentclass[11pt]{article}
\title{Policy Before Capacity: A Conditional Design and Evaluation Plan for Sovereign University AI Access}
\date{}
\author{SIGMA AI-assisted research synthesis}
\usepackage[T1]{fontenc}
\usepackage[utf8]{inputenc}
\usepackage[margin=1in]{geometry}
\usepackage{amsmath,amssymb,array,booktabs,longtable,pdflscape}
\usepackage{microtype}
\usepackage{titlesec}
\usepackage{tikz}
\usetikzlibrary{fit}
\emergencystretch=8em
\raggedbottom
\linespread{1.07}
\usepackage{xcolor}
\definecolor{blue}{HTML}{24485E}
\titleformat{\section}{\Large\bfseries\sffamily\color{blue}}{\thesection}{0.8em}{}
\titleformat{\subsection}{\large\bfseries\sffamily}{\thesubsection}{0.7em}{}
\setlength{\parskip}{.6em}
\setlength{\parindent}{0pt}
\begin{document}
\maketitle
\begin{center}\fbox{\parbox{.92\linewidth}{\textbf{Needs review: TeX verification is incomplete.} Review the archived proofreading diagnostics before relying on the typesetting.}}\end{center}
\begin{center}\small\textit{AI-generated research synthesis \textperiodcentered{} SIGMA}\end{center}

\begin{abstract}
Can universities share uneven AI capacity without surrendering data authority or exceeding hard budgets? This paper develops a bounded analytical comparison of campus-local inference, approved frontier APIs, a campus gateway with separately approved external access, and a hypothetical inter-campus federation. It treats useful research access as authorized parent-task completion that meets quality, provenance and latency requirements, rather than as model-call throughput. The analysis identifies a partition-related authorization limitation, conditional reservation and dispatch mechanisms, and counterexamples to unsafe liability release and aggregate quality reporting.

It also establishes a conditional economic threshold: federation lowers cost per useful task only when its proportional completion gain exceeds its proportional added cost. A repaired hypothetical workload illustrates why sensitive-task performance needs separate assumptions and evaluation. No deployment advantage is demonstrated. The resulting contribution is a policy-first architecture candidate and a proposed 90-day evaluation whose efficacy, safety, affordability and institutional approval gates remain distinct.
\end{abstract}
\section*{About this paper}
This paper was generated by SIGMA, an AI-assisted research system, from the project's recorded investigations, findings, and supporting evidence.
The analysis and conclusions reflect an automated synthesis of that research. Although SIGMA incorporates evidence tracking, critical review, and consistency checks, the resulting paper has not necessarily undergone independent verification or peer review.
Readers should evaluate its arguments, supporting evidence, and stated limitations before relying on its conclusions.
Research style: Distinguish algorithmic guarantees from tests, benchmarks and conjectures. Document counterexamples, complexity assumptions and reproducibility.
\section*{Evidence coverage limitation}
The evidence index omitted 9718 material saved research field fragment(s) across 10 source(s). These fragments were not reviewed for this paper. Interpretations that could depend on omitted material remain provisional; the complete frozen source snapshot and omission ledger are retained with the generation record.
\textit{Claims discussed: }C11, C13, C14, C16, C19
\section{Introduction and problem statement}
The relevant comparison is not simply local models against frontier models. Campus-local inference, approved API access, a campus gateway with separately approved external access, and a hypothetical federation expose different policy, capacity and administrative constraints. Federation must therefore demonstrate an advantage over the gateway baseline, not merely over an underprovisioned local service.

Research access requires both lawful processing and useful results. The comparison counts offered parent research tasks, including tasks that remain unserved, and distinguishes policy eligibility, admission and successful completion. Child model calls, agent actions, retries and reproducibility repeats consume resources but do not create additional parent-task successes.

The institutional requirement is continued home-campus authority over identity, classification, approved destinations, retention, auditing, contracts and quotas. Sensitive and derived content must remain within approved domains by default; ambiguous classification, service failure and budget exhaustion cannot create permission to use another destination. These are design requirements, not assertions about an implemented federation or established university policy.

\textit{Claims discussed: }C15, C19
\section{Methods and evidence}
The analysis used four explicitly defined service alternatives and a common parent-task accounting framework. It separated authorized destination eligibility from technical admission and joint quality, provenance and deadline success. A hypothetical workload worksheet supported conditional comparisons, while a stratum-specific counterexample exposed an invalid inference from aggregate local usefulness to sensitive-task performance.

Authorization was examined through paired executions that are indistinguishable to a disconnected admission point. Budget and dispatch safety were examined through proposed guarded state transitions, conditional invariant arguments and hand-worked adversarial schedules. These analyses were specifications and deductions, not executed fault tests, implementation inspections or formal verification.

Economic analysis used success-adjusted cost ratios and symbolic reservation limits rather than invented vendor rates. Pilot planning distinguished a proposed confidence-bound decision rule from approximate precision calculations and a later conservative analytical sample-size certificate. No empirical four-baseline evaluation, current-price retrieval or completed pilot is represented as having occurred.

\textit{Claims discussed: }C11, C12, C13, C14, C16, C17, C18
\section{Four service alternatives and the unit of useful access}
The comparison separates the scope of available destinations from the unit of useful research access: adding an execution destination does not itself authorize its use or establish that more research tasks can be completed. Four analytical alternatives distinguish campus-local service, approved external service, their coexistence, and a proposed inter-campus extension. These are comparison definitions, not observations of four deployed systems. In particular, a campus gateway with separately approved external access must not be conflated with either integrated frontier routing or a cooperative federation [C2, C3].

The supplied characterization of MindRouter provides the campus-gateway starting point. It attributes protocol translation across OpenAI, Ollama and Anthropic clients and Ollama/vLLM backends, WDRR fair-share scheduling, per-user quotas, audit logging and GPU telemetry to the published gateway. The starting references are the documentation at https://mindrouter.ai/documentation.html and the PEARC\textquoteright{}26 short paper at https://ai4ra.uidaho.edu/wp-content/uploads/2026/04/Sheneman\_MindRouter\_PEARC26\_Short\_Paper.pdf.

Their reported historical review date is 2026-10-31, but neither reference was retrieved for this analysis, so no actual access date, current contents or changed feature claims are established here; the capability list is a supplied description of published claims, not an independently measured result [C1].

The accompanying table fixes the alternatives by destination scope, treatment of campus-only demand and institutional boundaries. Separately approved external access means that external use requires its own approval; it does not imply that the campus gateway automatically selects or invokes a frontier provider. The federation alternative adds possible partner-campus destinations to the gateway baseline without relaxing the home campus\textquoteright{}s eligibility rules. Neither federation nor integrated frontier routing is attributed to existing MindRouter functionality [C1, C2].

The common outcome is a useful parent-task completion, rather than a successful model invocation. A parent task counts only if its processing is authorized and its result satisfies the prespecified quality rubric, required provenance and delivery deadline. These conditions define a joint event: a high-quality but unauthorized result, an inadequately documented result where provenance is required, or a result delivered after its deadline does not qualify. Child calls, tool invocations and retries consume resources within that parent task rather than creating additional successful tasks [C3].

Formally, fix a common offered cohort of $D$ parent tasks, its task mix and an evaluation horizon. For baseline $j$, define $y_{ij}=1$ when parent task $i$ satisfies all useful-completion requirements, and $y_{ij}=0$ otherwise; useful completions are $S_j=\sum_{i=1}^{D}y_{ij}$. Overall completion coverage is therefore $S_j/D$ for $D>0$, with offered but unserved tasks retained in the denominator. This accounting prevents a multi-call workflow from appearing more productive merely because it generates more invocations or retries [C3].

Coverage and model-quality comparisons require different denominators. Overall coverage includes demand that an alternative cannot serve, including campus-only tasks under the API-only alternative, whereas model-quality comparisons require matched policy-eligible task strata. Otherwise, excluding restricted work could make external service appear favorable without answering whether it provides access to the university\textquoteright{}s offered workload. A correct policy refusal is reported separately as a safety outcome, not counted as delivered research [C2, C3].

These definitions make the gateway-to-federation comparison a question about additional authorized, timely, rubric-passing and provenance-complete parent tasks. They do not permit additional destinations, nominal capacity or call volume to stand in for useful access. The resulting burden for the federation proposal is to improve task-level service within unchanged campus permissions, rather than to treat connectivity itself as an access gain [C2, C3].

\begin{landscape}
\small
\begin{longtable}{@{}>{\raggedright\arraybackslash}p{\dimexpr(\linewidth-2\tabcolsep*4)/210*42\relax}>{\raggedright\arraybackslash}p{\dimexpr(\linewidth-2\tabcolsep*4)/210*42\relax}>{\raggedright\arraybackslash}p{\dimexpr(\linewidth-2\tabcolsep*4)/210*42\relax}>{\raggedright\arraybackslash}p{\dimexpr(\linewidth-2\tabcolsep*4)/210*42\relax}>{\raggedright\arraybackslash}p{\dimexpr(\linewidth-2\tabcolsep*4)/210*42\relax}@{}}
\caption{Analytical service alternatives. Destination scope specifies where processing may be considered, not guaranteed admission, capacity or useful completion. No performance measurements are assigned to these alternatives.}\\
\toprule
\textbf{Alternative} & \textbf{Destination scope} & \textbf{Treatment of campus-only demand} & \textbf{Institutional boundaries} & \textbf{Evidence status}\\
\midrule
\endfirsthead
\toprule
\textbf{Alternative} & \textbf{Destination scope} & \textbf{Treatment of campus-only demand} & \textbf{Institutional boundaries} & \textbf{Evidence status}\\
\midrule
\endhead
\midrule
\multicolumn{5}{r}{\footnotesize Continued on next page}\\
\endfoot
\bottomrule
\endlastfoot
$\mathsf{L}$: campus-local inference only & Home-campus inference resources only & Considered locally; eligibility does not guarantee service & Home-campus processing domain; no external inference destination & Analytical baseline, not an empirically evaluated configuration here\\
$\mathsf{A}$: approved frontier API only & Separately approved frontier API destinations only & Unserved and retained in the overall coverage denominator & Home-campus approval and an external provider boundary & Analytical baseline, not an empirically evaluated configuration here\\
$\mathsf{G}$: campus gateway with separately approved external access & Campus inference plus external access approved separately; no assumed integrated frontier routing & Remains on the campus-local path & Campus processing domain plus separately approved provider boundaries & Gateway capabilities are supplied published-claim descriptions, not independently verified capabilities\\
$\mathsf{F}$: proposed cooperative federation added to $\mathsf{G}$ & Gateway baseline plus explicitly approved partner-campus destinations & Remains campus-only; federation grants no default permission for partner processing & Gateway-baseline boundaries plus partner-campus authority and agreement boundaries & Hypothetical extension, not an existing MindRouter feature or measured deployment\\
\end{longtable}
\par\smallskip{\footnotesize Evidence: [Project overview], [Research session 1, artifact 2], [Research session 1, artifact 5], [Research session 1, artifact 8]; claims: C1, C2, C3}
\end{landscape}
\textit{Claims discussed: }C1, C2, C3
\section{Home-campus authority before destination ranking}
Home-campus authorization must determine which destinations may be considered before capacity, quality or price can influence routing. The proposed rule therefore separates eligibility from ranking: it first constructs a campus-authorized destination set for each operation and only then compares quality estimates, deadline feasibility and incremental cost within that set. This ordering makes policy a constraint on optimization rather than a preference that a sufficiently attractive destination can outweigh [C4].

Eligibility must bind the authenticated principal and project to the operation\textquoteright{}s data classification, approved processing domains, retention conditions, audit requirements, contractual permissions, model and context requirements, and available quota. An unresolved classification produces an empty eligible set, not a low-confidence permission. If no destination qualifies, the permitted dispositions are local queuing where authorized, rejection, or a request for explicit approval; neither resource availability nor an existing financial reservation independently confers authority to disclose data [C4, C6].

The proposed classification boundary covers restricted, student, patient, controlled, sponsor-restricted and unpublished data, including attachments and derived prompts, outputs, logs, embeddings and caches. These artifacts remain within approved domains by default; classification cannot be limited to the original uploaded file while leaving subsequent representations unrestricted. Derived content should conservatively inherit the applicable restrictions of its inputs unless an authorized transformation or review changes that classification, while unknown or incompatible restrictions stop processing at an approved boundary.

Classification inheritance is a proposed institutional rule whose acceptance remains unresolved, not an established policy shared by participating campuses [C4, C19].

Enforcement must consequently extend beyond selecting a model endpoint. The proposed boundary register identifies the transfers and secondary stores that require authorization, including paths created by tools and backend services. Authorization must precede sensitive payload transfer: a partner\textquoteright{}s rejection after receiving the prompt cannot prevent the disclosure that has already occurred. Campus-local placement likewise does not, by itself, establish acceptable access, retention or logging [C4].

A minimal enforcement design would retain the home-campus identity and policy authority, local gateway and backends, protected audit storage, and controlled external-egress adapter. Direct tenant or backend paths that bypass the adapter would have to be denied, and protocol normalization would have to preserve policy-relevant attachments, tool instructions and metadata or reject unsupported semantics. Every agent or tool call and every destination change would require renewed eligibility checking against its actual content and destination. These requirements depend on complete mediation and trustworthy classification bindings; a routing recommendation alone establishes neither containment nor compliance by an authorized recipient [C4, C6].

Failure does not enlarge the eligible set. Outage, overload, provider rejection, timeout, retry, cancellation, exhausted budget or a better destination score must not silently authorize a previously prohibited route. A fallback is therefore a new candidate operation subject to renewed checks, rather than an exception to the original decision. This rule also prevents an agent\textquoteright{}s initially approved workflow from treating approval of its first model call as blanket permission for subsequent external tools [C4].

Policy freshness introduces a separate limitation that cannot be resolved by destination scoring. In the conditional partition model, let $H$ be the home-campus authority, $E$ a disconnected admission point, and $D$ a destination previously approved by $H$. The new request has no previously committed, request-specific authorization exempting it from subsequent revocation, and no alternative policy, credential or client channel communicates later policy changes to $E$. Immediate revocation means denying new authorization decisions ordered after the revocation takes effect; partition availability means admitting an otherwise eligible new request while disconnected when approval remains unchanged [C5].

Consider two executions with the same initial approval and partition. In execution $A$, $H$ leaves approval of $D$ unchanged; in execution $B$, $H$ revokes $D$ before the new request reaches $E$. The request, credentials, cached policy, clock readings and local state visible to $E$ are identical in both executions. A signature on the cached policy authenticates its historical issuer and contents but supplies no information distinguishing unchanged approval from an unseen revocation [C5].

A deterministic admission rule must therefore make the same decision in both executions. Admitting in $A$ to provide partition availability also admits in $B$, violating immediate revocation; reliably denying in $B$ also denies in $A$. Randomization does not remove the limitation: coupling the random choices yields \$\textbackslash{}Pr(\textbackslash{}mathrm\{admit\}\textbackslash{}mid A)=\textbackslash{}Pr(\textbackslash{}mathrm\{admit\}\textbackslash{}mid B)\$, so zero unauthorized-admission probability in $B$ requires zero admission probability in $A$. The proof candidate is conditional and not formally verified; it constrains the stated disconnected-authorization model rather than demonstrating a defect in MindRouter or an impossibility for every federation design [C5].

The design choices are consequently to deny or defer the affected requests, regain authoritative contact, or adopt explicitly campus-approved weaker revocation semantics. An unexpired lease changes the requirement only if the campus agrees that revocation may wait until lease expiry; a time-to-live alone does not provide immediate enforcement. A connected home-controlled decision service avoids the disconnected-authorizer premise, but its authorization ordering still does not recall previously disclosed data or guarantee that already authorized bytes cannot arrive after a later revocation [C5].

The proposed admission machine makes that ordering precise by treating preliminary approval as advisory and placing final authorization and consumption of the dispatch right together at the authoritative commit point. Under its stated assumptions, an applicable revocation committed before the dispatch grant prevents that grant, whereas a grant committed first can still result in later transmission or execution. The conditional guarantee is current-policy authorization at the decision point, not prevention of every post-revocation delivery. A campus requiring stronger delivery, execution or retention semantics would need additional recipient cooperation and explicit treatment of in-flight work [C5, C6].

Campus acceptance of classification inheritance, decision-time revocation and recipient obligations remains an input to feasibility; the supplied evidence establishes neither those approvals nor implemented enforcement. The resulting change to the capacity-sharing argument is substantive: another university\textquoteright{}s available GPU is a candidate resource only after the home campus authorizes the complete processing path, and cached approval cannot preserve both immediate revocation and new-request availability during the stipulated partition. Federation must therefore demonstrate benefit within those constraints, rather than count policy-prohibited routing or stale-policy admission as improved access [C4, C5, C19].

\begin{longtable}{@{}>{\raggedright\arraybackslash}p{\dimexpr(\linewidth-2\tabcolsep*2)/126*42\relax}>{\raggedright\arraybackslash}p{\dimexpr(\linewidth-2\tabcolsep*2)/126*42\relax}>{\raggedright\arraybackslash}p{\dimexpr(\linewidth-2\tabcolsep*2)/126*42\relax}@{}}
\caption{Proposed authorization boundary register. Entries specify required controls, not verified deployment properties.}\\
\toprule
\textbf{Boundary or artifact path} & \textbf{Required authorization scope} & \textbf{Disposition when unresolved}\\
\midrule
\endfirsthead
\toprule
\textbf{Boundary or artifact path} & \textbf{Required authorization scope} & \textbf{Disposition when unresolved}\\
\midrule
\endhead
\midrule
\multicolumn{3}{r}{\footnotesize Continued on next page}\\
\endfoot
\bottomrule
\endlastfoot
Client through gateway to local backend & Authenticated principal, purpose, classification, permitted access, retention and audit rules & Stop processing at an approved boundary; local placement is not sufficient permission\\
Gateway through external adapter to partner campus or provider & Home-campus approval of destination, content, contract, retention and audit conditions before payload transfer & Deny or defer external disclosure; remote rejection after transfer is too late\\
Backend or agent to a tool & Renewed eligibility for the individual call, its destination and its derived content & Do not inherit blanket permission from the parent workflow\\
Outputs, logs, embeddings, caches, storage and telemetry & Applicable inherited restrictions, approved storage or processing domain, access and retention controls & Keep artifacts inside approved domains; ambiguous classification fails closed\\
\end{longtable}
\par\smallskip{\footnotesize Evidence: [Research session 1, artifact 1], [Research session 1, artifact 3], [Project overview]; claims: C4}
\textit{Claims discussed: }C4, C5, C6, C19
\section{Reservations, one-use dispatch and conservative recovery}
The proposed admission machine makes policy authority and unresolved financial liability prerequisites for dispatch, rather than consequences to be checked after routing. Its contribution is a conditional argument: authorization must hold when the home authority commits a dispatch grant, and every potentially billable attempt must remain covered until its liability is resolved. These controls are proposed extensions, not established MindRouter capabilities; they address authorization and allocation safety, not task quality, availability or a demonstrated advantage from federation. [C6, C19]

The argument assumes complete mediation of billable actions, trustworthy classification bindings, a durable and linearizable home authority, atomic accounting across all applicable budget scopes, enforceable charge ceilings and authoritative settlement. Complete mediation includes retries, agent and tool subcalls, and any separately approved API account whose expenditure is claimed to fall under the same cap. Alternate credentials or hidden paid retries would place expenditure outside the invariant's domain. Neither implementation inspection nor executed fault tests establish these premises in the supplied evidence, and institutional acceptance of the required authorization and billing semantics remains unresolved. [C6, C19]

The request identity must distinguish a parent task from its operations and potentially billable attempts. Let $R$ identify the parent task, $O$ an immutable logical inference or tool operation, and $A$ one attempt identified by the home issuer, operation and server-issued sequence. An explicitly authorized retry receives a new attempt identity under the unchanged operation; a changed payload or destination requires a new operation linked to the same parent task. A client idempotency key, scoped to principal and project, returns the existing operation and status for the same manifest, while a conflicting manifest is rejected; neither an identifier nor a status response conveys sending authority. [C6]

The immutable manifest binds what will actually be transmitted, not merely the client's initial request. Protocol translation and attachment resolution therefore precede freezing, which binds the translated wire request, attachment bytes or immutable approved-domain references, destination, model request, resource ceilings, tool permissions and retention requirements. Unknown fields or mutable content invalidate validation. A protected digest can assist binding, but cannot establish that classification is correct; manifests, digests and deduplication records must themselves remain access-controlled within approved domains. [C6, C8]

Preliminary approval is advisory because policy, identity or cancellation state can change before dispatch. The proposed home-campus transaction authority orders those changes with attempt transitions and ledger updates, and the final dispatch-grant commit both checks current authorization and consumes the dispatch right. An applicable revocation committed first prevents the grant; a grant committed first authorizes only that immutable attempt. This is decision-time authorization: the designated sender may transmit after a later revocation, so the claim does not imply post-revocation non-delivery or recipient compliance. [C6, C8]

For accounting, let $s$ range over the applicable user, project, department and campus scopes. Within an accounting window, $B_s$ is the cap, $L_s$ settled cumulative liability, and $H_s$ the sum of reservations covering unresolved billable actions; the required initial condition is $L_s+H_s\leq B_s$ for every scope. Before dispatch, an attempt receives an enforceable maximum liability $r$, and reservation succeeds only if $L_s+H_s+r\leq B_s$ holds in every applicable scope. The transaction adds the hold to all scopes atomically or to none: the same liability is constrained at several organizational levels, not billed several times, and concurrent admissions cannot independently consume the same remaining allocation. [C6]

The bound must cover the attempt's total possible liability, including applicable ancillary charges and currency or tax exposure, rather than only an expected token charge. Without a finite enforceable bound, a paid action cannot be admitted while guaranteeing a finite cap against every allowed billing outcome. GPU service quotas and concurrency limits remain separate resources requiring their own admission controls; monetary headroom does not establish resource availability. [C6]

Conditional preservation follows by induction over ledger transitions. Admission preserves the inequality through its atomic guard, while final settlement replaces the reservation $r$ with actual liability $z$, where $0\leq z\leq r$, changing the scoped total to $L_s+H_s-r+z\leq L_s+H_s$. A proven unbilled rejection permits $z=0$; an ambiguous timeout, lost acknowledgement, cancellation request or WAN failure leaves the reservation unchanged. Durable attempt and settlement-event identities are needed to prevent repeated processing from settling or releasing the same hold twice. [C6]

The release rule is necessary, not merely conservative bookkeeping. In a hypothetical ledger with a cap of $10$, an attempt reserves $6$ and then times out without resolving its charge; releasing that hold permits a second liability of $6$, leaving possible total liability of $12$. Cancellation alone has the same defect unless it establishes final billing status. Retaining the first hold blocks the second admission under that cap, exchanging availability for continued coverage of the unresolved obligation; this hand-worked counterexample refutes premature release, not the conditional reservation invariant. [C7]

A durable grant-consumption record does not by itself prevent replay at the sender. The proposed burn-before-send adapter therefore consumes the grant before transmission and binds the acknowledgement to one original live continuation, protected by a one-shot transport latch. Sending authority is not reconstructed from a status query, queue, restart or takeover. Under these unimplemented sender assumptions, a lost acknowledgement or crash can sacrifice an otherwise authorized attempt, but cannot justify initiating it again under the same identity. [C8]

The resulting claim is at-most-one adapter invocation initiation per attempt, not exactly-once provider execution. It neither couples the database commit atomically to remote execution nor turns an uncertain response into evidence of non-execution. Any further attempt must obtain fresh authorization and financial coverage while the earlier unresolved hold remains in place. The accompanying fault schedules specify expected rejection or retention behavior, rather than observed test results. [C6, C8]

Settlement finality must also cover a still-possible delayed send. A current zero-charge report, local silence or the end of an invoice period does not establish that an original live continuation can no longer generate liability. Release therefore requires terminalization of that possibility or an enforceable final-liability guarantee covering delayed requests; an operational predicate establishing this condition remains open. Consequently, uncertainty holds may reduce availability indefinitely rather than safely expiring on a local timer. [C6, C8, C19]

Recovery must preserve financial exposure as well as reject replay. With a cap of $10$, suppose an old committed attempt can still incur $6$, but a stale backup omits its reservation; issuing fresh credentials and a new fenced issuer namespace can prevent reuse of old sending authority without restoring the missing obligation. Admitting a new liability of $6$ then permits exposure of $12$ despite successful replay prevention. New admissions must therefore wait until prior obligations are preserved, authoritatively reconciled or conservatively reserved against available funds; namespace isolation alone is an unsafe recovery shortcut. [C9]

The same liability-preservation requirement applies across accounting windows: unresolved obligations must remain against their original allocation or move to an explicitly funded liability account, rather than disappear at month-end. Cap reductions likewise cannot invalidate existing obligations, and allocation safety does not imply invoice-month cash-flow safety unless billing timing is also bounded. These results narrow the central argument: a hard budget requires bounded, continuously covered liability and an explicitly availability-sacrificing dispatch discipline, not merely a quota check before routing. Neither property requires federation, and neither demonstrates that federation supplies useful additional access. [C6, C19]

\begin{quote}
\textbf{Diagram (readable relationship list).} Proposed policy-first control and data paths for a campus gateway and hypothetical external extensions. Every destination requires home-campus approval; the diagram does not assert existing MindRouter federation features or verified enforcement. Uncertain transmission or billing retains liability rather than enabling an unapproved fallback.
\begin{description}
\item[\textbf{Node 1.}] Research client and tenant input
\item[\textbf{Node 2.}] Identity, classification and eligible-destination check
\item[\textbf{Node 3.}] Authoritative authorization and atomic scoped reservation commit
\item[\textbf{Node 4.}] Proposed one-use, burn-before-send adapter
\item[\textbf{Node 5.}] Approved campus-local backend
\item[\textbf{Node 6.}] Hypothetical separately approved peer-campus destination
\item[\textbf{Node 7.}] Separately approved frontier workflow
\item[\textbf{Node 8.}] Durable audit, unresolved holds and authoritative reconciliation
\item[\textbf{Edge 1.}] Node 1 points to node 2: Treat requests and derived content as policy-bearing input
\item[\textbf{Edge 2.}] Node 2 points to node 3: Only an eligible, immutable operation may proceed
\item[\textbf{Edge 3.}] Node 3 points to node 4: Consume the authoritative dispatch right
\item[\textbf{Edge 4.}] Node 4 points to node 5: Authorized local execution
\item[\textbf{Edge 5.}] Node 4 points to node 6: Only if the home campus separately approves the peer
\item[\textbf{Edge 6.}] Node 4 points to node 7: Only if external access and liability are approved
\item[\textbf{Edge 7.}] Node 3 points to node 8: Persist scoped commitments and authorization evidence
\item[\textbf{Edge 8.}] Node 4 points to node 8: Record observations; ambiguity does not release liability
\item[\textbf{Edge 9.}] Node 8 points to node 3: Outstanding liabilities reduce admission headroom
\item[\textbf{Boundary 1.}] Untrusted tenant input: node 1
\item[\textbf{Boundary 2.}] Proposed home-campus control and approved local processing domain: node 2, node 3, node 4, node 5, node 8
\item[\textbf{Boundary 3.}] Distinct peer-campus authority and processing domain: node 6
\item[\textbf{Boundary 4.}] Distinct provider contractual and processing domain: node 7
\end{description}
\par\smallskip{\footnotesize Evidence: [Project overview], [Research session 1, artifact 3], [Research session 1, artifact 4], [Research session 1, artifact 7]; claims: C4, C6, C7, C8, C9}
\end{quote}
\begin{longtable}{@{}>{\raggedright\arraybackslash}p{\dimexpr(\linewidth-2\tabcolsep*2)/126*42\relax}>{\raggedright\arraybackslash}p{\dimexpr(\linewidth-2\tabcolsep*2)/126*42\relax}>{\raggedright\arraybackslash}p{\dimexpr(\linewidth-2\tabcolsep*2)/126*42\relax}@{}}
\caption{Discriminating fault schedules for the proposed admission and dispatch mechanisms. Outcomes are analytical expectations and proposed acceptance criteria, not executed tests.}\\
\toprule
\textbf{Adversarial schedule} & \textbf{Required behavior under the proposal} & \textbf{Discriminating failure}\\
\midrule
\endfirsthead
\toprule
\textbf{Adversarial schedule} & \textbf{Required behavior under the proposal} & \textbf{Discriminating failure}\\
\midrule
\endhead
\midrule
\multicolumn{3}{r}{\footnotesize Continued on next page}\\
\endfoot
\bottomrule
\endlastfoot
An applicable revocation commits after preliminary approval but before the dispatch-grant commit. & The authoritative commit rejects the grant despite the earlier approval or existing reservation. & A grant is committed using the superseded authorization.\\
Concurrent attempts contend for remaining allocation that cannot cover both reservations. & The multi-scope transactions serialize; only reservations satisfying every applicable cap commit. & Both attempts consume the same headroom, or an attempt receives only a partial set of scoped holds.\\
A grant is consumed, its acknowledgement is lost, and a status query, restart or takeover follows. & The consumed grant does not yield reconstructed sending authority; financial coverage persists while liability remains possible. & Recovery or a status response initiates a replacement invocation under the same attempt.\\
The original continuation encounters a duplicated acknowledgement or repeated send trigger. & The acknowledgement remains bound to that continuation, and its one-shot latch permits at most one invocation initiation. & More than one adapter invocation initiation occurs for the attempt.\\
The same authoritative settlement event is delivered more than once. & Durable settlement-event identity prevents repeated replacement or release of the reservation. & Duplicate settlement reduces held liability or changes settled liability a second time.\\
\end{longtable}
\par\smallskip{\footnotesize Evidence: [Research session 1, artifact 4], [Research session 1, artifact 7]; claims: C6, C8}
\textit{Claims discussed: }C6, C7, C8, C9, C19
\section{Capacity, fairness and reproducibility are separate service obligations}
Capacity sharing supports the paper\textquoteright{}s central argument only if it improves authorized, timely research outcomes; scheduling weights, GPU utilization and reproducible execution are separate obligations rather than interchangeable evidence of useful access. In particular, weighted deficit round-robin scheduling (WDRR) alone establishes neither unconditional latency bounds nor starvation freedom. The service comparison must therefore distinguish allocation policy from the compatible capacity, workload conditions and execution records needed to deliver a parent task successfully. [C3, C10]

Capacity should be inventoried by the work it can execute, not merely by aggregate GPU count. The proposed comparison must record heterogeneous GPU memory, model and context fit, memory headroom, available compatible capacity and admitted demand. An idle device that cannot accommodate a task\textquoteright{}s required model does not resolve that task\textquoteright{}s capacity shortage, and permission to use a remote campus does not establish that a compatible device is available there. This distinction makes federation\textquoteright{}s contribution conditional on usable capacity within the task\textquoteright{}s approved destinations, rather than on the nominal size of the combined fleet. [C10, C19]

The scheduling limitation is substantive, not simply an absence of benchmark results. A positive scheduling weight cannot supply service when compatible capacity is unavailable; likewise, weights alone do not bound delays caused by blocking, queues, cold starts or maintenance. Peak demand and interactive-versus-batch contention must therefore be explicit test conditions, rather than circumstances omitted from an otherwise favorable latency report. These constraints rule out an unconditional WDRR service guarantee, but do not rule out useful service under measured capacity and demand conditions. [C10]

The proposed evidence register separates the service obligations that the pilot must assess across the four alternatives. It is a specification for measurement, not a report of achieved performance: no executed scheduling benchmark, campus-specific capacity study or measured federation advantage is supplied. In particular, the fairness reference allocation and its observation windows still require institutional approval, so a favorable allocation result cannot yet be defined solely by observing scheduler weights. [C10, C16, C19]

Quality assessment should operate on the delivered artifact against a rubric fixed before execution. The proposed rubrics require registered tests and critical-correctness checks for coding, domain-specific numerical checks for routine analysis, expert assessment of conclusions and evidence for scientific reasoning, and task-specific references for multimodal work. Agentic workflows must achieve their objective without unauthorized actions, while reproducible batches require a registered rerun and tolerance criterion. Two blinded assessors would score artifacts independently, with a third adjudicator resolving disagreement; agreement and ambiguous judgments should be reported rather than suppressed. [C3, C16]

Provenance recording and reproducibility testing must also remain distinct. The proposed provenance record includes the available model and version, configuration, prompts, tools and retrieval snapshots; a version label alone does not demonstrate that an artifact can be reproduced. Model and workflow versions should be locked where possible, changes recorded, and reruns assessed against the task\textquoteright{}s prespecified tolerance. Reproducibility repeats consume resources but do not become additional parent-task successes. [C3, C16]

Latency should run from submission to the assigned workflow through delivery of the completed artifact, including approval handling, queueing, cold starts, tools, retries and recovery. First-token latency is therefore insufficient for the joint completion endpoint, and later human grading does not extend the delivery interval. Class-specific deadlines must be approved before enrollment; sensitive tasks inherit their functional task\textquoteright{}s deadline while retaining their stricter processing-domain requirements. Throughput should report useful parent-task completions over the observation period, with child calls and retries recorded separately as resource consumption. [C3, C16]

Availability and fairness reporting must retain the demand that the service does not deliver. Model-fit failures, queue rejection, budget exhaustion, provider rejection, timeouts and non-delivery by the deadline remain outcomes rather than post hoc exclusions. Across all four alternatives, requested, authorized and admitted demand should be reported by relevant task and policy stratum; the primary paired comparison must retain the same enrolled denominator in both arms. Correct policy refusal is a safety outcome, not a useful research completion, and tasks prohibited from paired evaluation must remain visible in coverage reporting without being exported for comparability. [C3, C16]

The fairness assessment should relate observed service to an approved reference allocation and observation window, while recording compatible capacity, blocking and scheduling quanta. Results should distinguish campus and workload strata, including interactive and batch service, so that aggregate utilization or completion does not substitute for evidence about prolonged waiting and unserved demand. Where a scheduling guarantee cannot be established, the comparison should explicitly report its absence alongside observed service. Neither favorable finite observations nor positive WDRR weights establish starvation freedom. [C10, C16]

Performance measurements also require comparable load conditions. The proposed gateway-versus-federation comparison holds local allocations, paid-access entitlements and parent-task spending limits constant, while declaring the additional remote capacity available to the federation. Isolated capacity slices with matched background arrivals, or separate matched replay windows, should prevent duplicate study executions from competing in one queue and then being presented as ordinary-service latency.

Peak, cold-start and maintenance conditions must be included, versions and changes recorded, and answer sharing or shared cross-arm caches prevented; controlled replay evidence should remain distinguishable from live-service observations affected by shared queues and outages. [C10, C16]

These distinctions change what federation must demonstrate: more nominal capacity or higher call throughput is not enough. Its added capacity must produce useful parent-task completions under the required quality, provenance and deadline rules, while service and fairness obligations remain independently satisfied. The simpler campus gateway with separately approved external access therefore remains the relevant comparator until the federation\textquoteright{}s additional resources yield a measured service benefit sufficient to support the separate efficacy and administrative-cost gates. [C3, C10, C16, C19]

\begin{longtable}{@{}>{\raggedright\arraybackslash}p{\dimexpr(\linewidth-2\tabcolsep*2)/118*34\relax}>{\raggedright\arraybackslash}p{\dimexpr(\linewidth-2\tabcolsep*2)/118*42\relax}>{\raggedright\arraybackslash}p{\dimexpr(\linewidth-2\tabcolsep*2)/118*42\relax}@{}}
\caption{Proposed service-evidence register. Each row identifies a separate measurement obligation; none represents an achieved service level.}\\
\toprule
\textbf{Service obligation} & \textbf{Evidence to collect} & \textbf{Interpretive constraint}\\
\midrule
\endfirsthead
\toprule
\textbf{Service obligation} & \textbf{Evidence to collect} & \textbf{Interpretive constraint}\\
\midrule
\endhead
\midrule
\multicolumn{3}{r}{\footnotesize Continued on next page}\\
\endfoot
\bottomrule
\endlastfoot
Compatible capacity & GPU memory and headroom; model/context fit; available compatible capacity; admitted workload. & Aggregate GPU count or utilization does not establish capacity for a particular authorized task.\\
Responsiveness & Submission-to-delivery latency, including approval, queues, blocking, cold starts, tools, retries and recovery; declared peak and maintenance conditions. & First-token latency and warmed idle execution do not establish end-to-end deadline feasibility.\\
Fair allocation & Approved reference allocation and observation windows; scheduling weights and quanta; compatible capacity; waiting and unserved demand by campus and workload stratum. & WDRR weights alone establish neither unconditional latency bounds nor starvation freedom.\\
Useful throughput and availability & Requested, authorized and admitted parent tasks; joint quality/provenance/deadline successes; rejection, outage and non-delivery outcomes; child-call and retry consumption. & Unserved demand remains visible, and additional calls are not additional research successes.\\
Provenance and reproducibility & Available model/version, configuration, prompts, tools and retrieval snapshots; version changes; registered reruns and tolerance results. & A version label is not a reproducibility result, and reruns do not enlarge the parent-task success denominator.\\
\end{longtable}
\par\smallskip{\footnotesize Evidence: [Project overview], [Research session 1, artifact 3], [Research session 1, artifact 6], [Research session 1]; claims: C3, C10, C16, C19}
\textit{Claims discussed: }C3, C10, C16, C19
\section{Success-adjusted costs and the federation break-even condition}
Federation has an economic advantage only if its gain in useful parent-task completions outweighs its net added cost on a common accounting basis. The relevant comparison is therefore not GPU count or model-call throughput, but expenditure per authorized result that meets the task\textquoteright{}s quality, provenance and delivery requirements. This section establishes that conditional threshold and specifies the cost accounts needed to apply it without confusing existing investment, avoidable spending and long-run ownership cost (C3, C11, C12).

The comparison must hold the offered workload, task mix, evaluation horizon and acceptance rules fixed. For each service alternative, the useful-completion denominator counts successful parent tasks; child calls, retries and repeated Investigator/Skeptic exchanges consume resources without becoming additional successes. Offered but unserved tasks remain in completion-coverage denominators, while correct policy refusals are reported as safety outcomes rather than delivered research results. Cost per useful completion should therefore be reported alongside overall coverage, with matched policy-eligible strata used for model-quality comparisons so that excluding restricted work does not create an apparent efficiency advantage (C3).

Three accounting views answer different decisions and should remain separate. The sunk-cost account records prior expenditure; the incremental account identifies costs added or avoided by the proposed service change; and the fully loaded account includes ownership and operating resources over the chosen horizon. Acquisition and replacement cash flows must also be reconciled with depreciation rather than counted again as independent ownership expenses. The proposed accounting table identifies the required entries and allocation questions; it is an input specification, not a populated cost comparison (C12).

For local capacity, utilization belongs in both the operating account and the interpretation of productive capacity: idle resources must not disappear from ownership costs merely because they deliver no useful tasks. Power and cooling, storage and network, hardware support, procurement and staff operations must be included with explicit scopes that prevent overlapping allocations. For API use, dated input, output and cached-token rates must be applied to the complete parent-task workflow, including repeated calls, retries, tools, embeddings, egress and applicable commitments.

Failed workflows can consume resources without increasing useful completions, so pricing a single successful inference is not an adequate substitute for workflow accounting (C3, C12).

The supplied evidence does not establish current API rates, campus cost inputs, opening liabilities, compatible capacity or active-researcher denominators sufficiently to produce a sourced four-alternative comparison. A hypothetical worksheet supports conditional accounting reasoning, but its complete numerical ledger is not available here, and reservation scenarios are neither dated token prices nor observed expenditure. Consequently, campus-specific monthly and annual totals, numerical break-even points, cost per active researcher and cost per useful task cannot yet be reported. The accounting specification must be populated with non-overlapping campus records and dated contractual inputs before it supports a procurement decision (C12).

The federation threshold itself does not require those missing numerical inputs. Let $C_G$ and $S_G$ denote the gateway baseline\textquoteright{}s cost and useful completions for the fixed workload and horizon, and let federation change them by $\Delta C$ and $\Delta S$, respectively. Assume $C_G>0$, $S_G>0$, $\Delta C>0$ and $S_G+\Delta S>0$, using the same cost basis for both alternatives. Federation then lowers cost per useful completion exactly when the following inequality holds (C11):
$\frac{C_G+\Delta C}{S_G+\Delta S}<\frac{C_G}{S_G}.$

Both completion denominators are positive, so cross-multiplication preserves the inequality and gives $(C_G+\Delta C)S_G<C_G(S_G+\Delta S)$. Cancelling the common product yields $S_G\Delta C<C_G\Delta S$; division by the positive product $C_GS_G$ produces the proportional condition:
$\frac{\Delta S}{S_G}>\frac{\Delta C}{C_G}.$
Equality leaves unit cost unchanged. In particular, a positive completion gain smaller than the proportional cost increase raises, rather than lowers, cost per useful completion; increased useful throughput alone is therefore insufficient to establish an economic advantage (C11).

A limiting case makes the conditional nature of the result explicit. Suppose federation supplies zero usable additional peer capacity, preserves identical service opportunities, produces no offsetting benefit and adds strictly positive overhead on the selected cost basis. Useful completions are then unchanged while expenditure rises, so the gateway baseline delivers the same useful results at lower cost. Policy exclusions, incompatible capacity or inability to meet deadlines can make nominal peer resources unusable, but their incidence is not established by the supplied evidence. This is a counterexample to an unconditional federation advantage, not a general argument against cooperation (C15).

Applying the threshold requires sensitivity analysis of both sides, not merely varying an API rate while holding successful output constant. The proposed comparison must examine task mix, demand, utilization, workflow amplification, compatible capacity and net coordination overhead without extrapolating constant success fractions through saturation. It must also retain independent safety, mandatory-coverage and hard-budget gates: a lower cost ratio does not authorize otherwise prohibited processing or establish affordability.

The economic argument thus changes the burden of justification from demonstrating spare peer GPUs to demonstrating enough additional useful completions to exceed the proportional net cost increase relative to gateway service with separately approved external access (C3, C11, C12, C15).

\begin{landscape}
\small
\begin{longtable}{@{}>{\raggedright\arraybackslash}p{\dimexpr(\linewidth-2\tabcolsep*3)/168*42\relax}>{\raggedright\arraybackslash}p{\dimexpr(\linewidth-2\tabcolsep*3)/168*42\relax}>{\raggedright\arraybackslash}p{\dimexpr(\linewidth-2\tabcolsep*3)/168*42\relax}>{\raggedright\arraybackslash}p{\dimexpr(\linewidth-2\tabcolsep*3)/168*42\relax}@{}}
\caption{Required cost-accounting specification for a common-horizon comparison. Entries describe evidence and allocation rules to be established, not measured costs or verified prices (C12).}\\
\toprule
\textbf{Cost component} & \textbf{Sunk or opening position} & \textbf{Incremental account} & \textbf{Fully loaded account and reconciliation}\\
\midrule
\endfirsthead
\toprule
\textbf{Cost component} & \textbf{Sunk or opening position} & \textbf{Incremental account} & \textbf{Fully loaded account and reconciliation}\\
\midrule
\endhead
\midrule
\multicolumn{4}{r}{\footnotesize Continued on next page}\\
\endfoot
\bottomrule
\endlastfoot
GPU acquisition, depreciation and replacement & Record existing purchases, installed assets and replacement timing. & Identify additional acquisition or replacement spending attributable to the alternative; establish avoidability. & Allocate ownership cost over useful life and reconcile it with purchase and replacement cash flows without double counting.\\
Utilization and idle capacity & Inventory existing capacity and its compatibility with the workload. & Establish whether added capacity is usable for authorized, deadline-feasible work. & Retain idle and stranded capacity in ownership accounting; do not equate nominal GPU availability with useful completions.\\
Power and cooling & Document the existing operating scope and allocation. & Identify changes attributable to the proposed service. & Include power and cooling consistently with utilization and idle operation.\\
Storage and network & Record existing provision and its allocation. & Identify added storage, network and applicable egress requirements. & Use non-overlapping allocations across local operations, external services and federation.\\
Hardware support, procurement and staff operations & Document existing commitments and staff responsibilities. & Validate added or avoidable support, procurement and operating effort. & Include the complete service scope while preventing overlapping staff or support charges.\\
API inference and workflow charges & Establish existing commitments and unresolved opening liabilities. & Apply dated input, output and cached-token rates to all workflow calls; include retries, tools, embeddings, egress and applicable commitments. & Reconcile complete parent-task consumption with charges; do not substitute reservation ceilings for observed spending.\\
Federation coordination and offsetting benefits & Identify responsibilities and resources already assigned to gateway service. & Establish net added overhead and validate any claimed offsetting benefit. & Apply the same cost basis to gateway and federation accounts before calculating their cost difference.\\
\end{longtable}
\par\smallskip{\footnotesize Evidence: [Project overview], [Research session 1, artifact 2], [Research session 1, artifact 5], [Research session 1]; claims: C3, C11, C12, C15}
\end{landscape}
\textit{Claims discussed: }C3, C11, C12, C15
\section{A repaired stratified workload illustrates the decision boundary}
A stratified workload makes the policy-first decision boundary explicit: sufficient aggregate completion does not establish service for a mandatory campus-only stratum, and additional destinations do not necessarily add useful capacity. The following monthly worksheet compares the four service alternatives under declared admission and usefulness assumptions, rather than treating those assumptions as measured performance. Its proposed targets are $1400$ useful parent-task completions out of $2000$ offered tasks, or $70\%$, including at least $160$ useful completions among $200$ sensitive campus-only tasks. These are analytical targets, not established university policy. [C2, C13, C14]

The hypothetical cohort comprises $100$ distinct active researchers and $2000$ monthly parent tasks: $600$ coding tasks, $400$ routine analyses, $300$ demanding scientific-reasoning tasks, $200$ multimodal or long-context tasks, $200$ agentic workflows, $200$ campus-only sensitive analyses, and $100$ reproducible batches. These are disjoint accounting bins for this worksheet, not a claim that sensitivity constitutes a separate technical task type. Usefulness denotes the joint event of authorized completion satisfying the task's quality, deadline, and provenance requirements; calls, tool invocations, and retries remain descendants of one parent task rather than additional successes. [C13, C14]

Assume deadline-feasible local admission capacity of $1000$ tasks per month, allocated to all $200$ sensitive tasks and $800$ others. The illustrative nonsensitive allocation is $400$ coding tasks, $300$ routine analyses, and $100$ batches. This leaves $300$ locally suitable overflow tasks and $700$ tasks whose capability requirements are assumed unavailable locally, giving the gateway plan $1000$ distinct residual API candidates without double counting. The approved-API-only plan instead has $1800$ eligible tasks, because the $200$ sensitive tasks remain campus-only; federation does not itself change that restriction. These are fixed scenario assumptions about permissions, capability, and capacity. [C2, C14]

The first result is a counterexample to inferring sensitive coverage from aggregate usefulness. Let $q_{L,S}$ and $q_{L,N}$ denote average useful-completion probabilities among the sensitive and other local admissions. An aggregate local probability of $0.80$ establishes only $200q_{L,S}+800q_{L,N}=800$. Choosing $q_{L,S}=0$ and $q_{L,N}=1$ satisfies that equation while providing no expected sensitive completions; equivalently, exact counts of zero sensitive successes and $800$ other successes preserve the aggregate result of $800/1000$. Giving sensitive tasks admission priority cannot repair this inference, because admission is not useful completion. [C13]

The repaired scenario therefore adds separate premises: $q_{L,S}=q_{L,N}=0.80$. Linearity of expectation then gives $200(0.80)=160$ expected sensitive completions and $800(0.80)=640$ expected other local completions, without requiring independence. More generally, with $a_S$ sensitive admissions, the expected sensitive target requires $a_Sq_{L,S}\geq160$; with all $200$ admitted, it requires $q_{L,S}\geq0.80$, whereas fewer than $160$ admissions cannot suffice even with perfect usefulness. The sensitive-stratum probability is thus an additional substantive assumption, not information recovered from the aggregate. [C13, C14]

Expected counts must also remain distinct from realized service. A common shock under which all $200$ sensitive tasks succeed with probability $0.80$ and all fail with probability $0.20$ preserves each task's assumed marginal probability and the expected count of $160$, yet misses the sensitive target with probability $0.20$. The realized criterion would instead require observed counts $Y_S\geq160$ and $Y\geq1400$. The repaired expectation calculation supplies neither a probability of meeting both criteria nor a confidence guarantee. [C14]

For paid admission, let $B$ be the available monthly API reservation envelope and let $r>0$ be a common, enforceable maximum liability per admitted parent task. The worksheet assumes no opening liability or minimum commitment, nonbinding subsidiary quotas, sufficient API capacity, and no reuse of reservation savings during the monthly plan. Under these premises, the gateway admits $n=\min(1000,\lfloor B/r\rfloor)$ distinct API tasks. Their average useful-completion probability is assumed to remain $0.90$ for every selected cohort, including budget-truncated cohorts; a full-cohort aggregate probability alone would not justify that assumption. [C14]

The gateway's expected total is consequently $800+0.90n$, with sensitive completions remaining $160$ in expectation. Meeting the overall expected-count target requires $800+0.90n\geq1400$, hence $n\geq\lceil(1400-800)/0.90\rceil=667$ and, because the residual candidate limit is $1000$, equivalently $\lfloor B/r\rfloor\geq667$. The table distinguishes this conditional gateway result from the different failure modes of local-only and API-only service. It also records federation only in the scoped case where no additional usable peer service is available; this is a comparison of declared plans, not an optimization over every possible routing allocation. [C2, C14, C15]

For positive API envelopes of $2000$, $4000$, and $8000$ dollars, the expected-count threshold is respectively $r\leq2000/667$, $r\leq4000/667$, and $r\leq8000/667$ dollars per parent task\textemdash{}approximately $2.9985$, $5.9970$, and $11.9940$ dollars. These are allowable maximum reservation liabilities, not token prices, average task costs, or predicted spending. At $B=0$, none of the declared plans passes both expected-count targets; if $r$ is unknown or unenforceable, the rule does not justify paid admission even when $B>0$. Changes in local usefulness, selected API quality, residual demand, or feasible capacity require recomputing the boundary rather than carrying forward the threshold of $667$. [C14]

The numerical workload and researcher denominator are hypothetical inputs, not an observed campus demand profile. Dated API rates, enforceable charging bounds, campus cost inputs, opening liabilities, and compatible capacity have not been established sufficiently to convert this worksheet into campus-specific monthly or annual totals, break-even estimates, or costs per useful task. Its reservation thresholds therefore identify the conditions under which the declared plan passes expected-count gates, not an affordable procurement choice. [C12, C14]

The federation comparison has a similarly precise scope. If federation supplies zero usable additional peer capacity, preserves identical service opportunities and non-federation entitlements, provides no offsetting reliability, scheduling, or administrative benefit, and adds overhead $H>0$ on the chosen accounting basis, then $C_F=C_G+H>C_G$ while useful completions remain equal. Under those conditions, the gateway has lower cost for the same outcomes; zero additional GPU capacity alone would not establish dominance if federation improved another consequential service dimension.

The repaired workload thus replaces an aggregate feasibility assertion with two separate requirements: evidence for mandatory-stratum usefulness and a positive service contribution sufficient to justify federation's added cost. [C13, C14, C15]

\begin{landscape}
\small
\begin{longtable}{@{}>{\raggedright\arraybackslash}p{\dimexpr(\linewidth-2\tabcolsep*4)/173*42\relax}>{\raggedright\arraybackslash}p{\dimexpr(\linewidth-2\tabcolsep*4)/173*42\relax}>{\raggedleft\arraybackslash}p{\dimexpr(\linewidth-2\tabcolsep*4)/173*14\relax}>{\raggedright\arraybackslash}p{\dimexpr(\linewidth-2\tabcolsep*4)/173*33\relax}>{\raggedright\arraybackslash}p{\dimexpr(\linewidth-2\tabcolsep*4)/173*42\relax}@{}}
\caption{Conditional expected-count comparison for the repaired monthly workload. Define $m_A=\min(1800,\lfloor B/r\rfloor)$ and $n=\min(1000,\lfloor B/r\rfloor)$. The API usefulness premise applies to each selected cohort. Targets concern expectations, not guaranteed realized counts.}\\
\toprule
\textbf{Declared service plan} & \textbf{Admissions} & \textbf{Expected sensitive completions} & \textbf{Expected total useful completions} & \textbf{Relation to the two expected-count targets}\\
\midrule
\endfirsthead
\toprule
\textbf{Declared service plan} & \textbf{Admissions} & \textbf{Expected sensitive completions} & \textbf{Expected total useful completions} & \textbf{Relation to the two expected-count targets}\\
\midrule
\endhead
\midrule
\multicolumn{5}{r}{\footnotesize Continued on next page}\\
\endfoot
\bottomrule
\endlastfoot
$\mathsf{L}$ & $200$ sensitive and $800$ other local tasks & $160$ & $800$ & Meets the sensitive target under the repaired premise, but fails the overall target.\\
$\mathsf{A}$ & $m_A$ approved API tasks; no sensitive tasks & $0$ & $0.90m_A\leq1620$ & Fails the sensitive target at every API envelope, even when the overall target is met.\\
$\mathsf{G}$ & $1000$ local tasks plus $n$ distinct approved API tasks & $160$ & $800+0.90n$ & Meets both expected-count targets exactly when $\lfloor B/r\rfloor\geq667$.\\
$\mathsf{F}$, scoped no-extra-service case & Same local and API plan as $\mathsf{G}$; no additional usable peer admissions & $160$ & $800+0.90n$ & Has the same expected-count boundary as $\mathsf{G}$; positive net overhead without an offsetting benefit increases cost.\\
\end{longtable}
\par\smallskip{\footnotesize Evidence: [Research session 1, artifact 8]; claims: C2, C14, C15}
\end{landscape}
\textit{Claims discussed: }C2, C12, C13, C14, C15
\section{A gated 90-day pilot rather than a deployment claim}
The proposed 90-day pilot should test whether federation adds a material, policy-permitted completion benefit over gateway access, rather than treat shared capacity as a deployment recommendation. Its dependency order is source and institutional approval, executable rubrics and fault-test preparation, resource-feasibility certification, paired enrollment, and a fixed final analysis. The four baselines remain necessary for coverage and cost comparisons, but the primary efficacy contrast is the hypothetical federation against the campus gateway with separately approved external access. This separates the question of federation\textquoteright{}s incremental value from the broader comparison of local and external service alternatives (C16, C19).

The first gate requires verification of the MindRouter documentation and PEARC\textquoteright{}26 paper, with actual access dates and versions, before their capability descriptions become premises of the implementation plan. The reported review date of 2026-10-31 is historical information, not evidence of a new source inspection. Campus security, data owners, legal counsel, procurement and finance must also approve processing domains, classification inheritance, retention, revocation and in-flight semantics, contractual obligations, enforceable charge bounds and ownership of unresolved liabilities.

The supplied evidence contains neither this source-verified capability review nor institutional approvals, an inspected enforcement implementation, campus-specific cost and workload inputs, or a measured federation advantage; these are prerequisites to enrollment, not administrative details to resolve during real-data testing (C12, C19).

Before those approvals, preparation should use public or synthetic tasks, including synthetic sensitive-data fixtures. Faculty and student representatives should approve task-specific quality rubrics, provenance requirements, delivery deadlines and accessibility needs, while computing and administrative owners specify service, fairness and administrative-cost acceptance rules. Coding tests, known-result analysis checks, blinded scientific-reasoning rubrics, labeled multimodal fixtures and agent root-goal checks are candidate assessment methods, not completed instruments. Each must become executable before the population and enrollment target are frozen; a task category alone does not define a reproducible pass rule (C16, C19).

Fault-test preparation is a separate dependency from task-quality preparation. The proposed tests should cover stale policy and credentials, malicious-tenant bypass, ambiguous classification, WAN partitions, campus or provider outages, provider rejection, ambiguous charged timeouts, cancellation, concurrent quota claims and unsafe fallback. They should also exercise the two consequential recovery obligations retained by the design: restoring a ledger must preserve outstanding financial exposure, and releasing a reservation requires evidence that excludes liability from a still-possible delayed send.

These are proposed tests, not executed results; successful fault injection would provide bounded implementation evidence rather than a general safety proof (C19).

Resource certification must precede confirmatory enrollment because statistical adequacy does not establish affordability or schedulability. For each budget scope $s$ and accounting period $t$, the proposed plan requires $S_{st}+R_{st}+F_{st}+\sum_{i,b}r_{ibst}\leq B_{st}$, where the terms represent opening settled charges, outstanding reserved liabilities, other fixed obligations, enforceable maximum liabilities for each parent task and arm, and the approved cap. The envelopes must include permitted child calls, retries, tools, embeddings, egress and other charge components, including charges through separately approved accounts.

Finance must populate these quantities from dated prices and contractual bounds; neither hypothetical reservation scenarios nor the analytical sample-size certificate supply the missing four-baseline cost comparison (C12, C18).

Computing must independently certify a time-indexed execution plan that respects compatible model capacity, GPU memory, queue isolation, cold starts, maintenance and deadlines, together with sufficient operations and blinded-grading staff. Aggregate GPU-hours cannot establish that a particular model can serve a particular task before its deadline. WDRR weights likewise establish neither unconditional latency nor starvation freedom, so fairness must be assessed against an approved reference allocation with adequate exposure in compatible service pools. Insufficient exposure should leave fairness unassessed, rather than turn an absence of observed failures into a pass (C10, C19).

For the efficacy comparison, let $G$ denote the gateway with separately approved external access and $F$ the same service augmented by approved remote-campus execution. Both arms should use the same home-campus identity, classification, retention, audit and destination rules, local allocations, paid-access entitlements and parent-task spending limits. The additional remote capacity in $F$ is the intervention; its operating and agreement costs must therefore be counted rather than treated as free. Federation and frontier routing remain hypothetical extensions, and neither an outage nor a better expected result permits either arm to bypass destination approval or financial mediation (C16, C19).

The sampling unit is a parent-task test episode containing one execution in each arm. A parent task is one registered research objective with an input snapshot, requested output, frozen rubric, deadline, classification, provenance requirements and total resource envelope; agent calls, tools, retries and reproducibility repeats are descendants, not additional observations. Define $Y_b=1$ only when arm $b$ satisfies the quality rubric, authorized processing, required provenance and delivery deadline jointly, and set $Y_b=0$ otherwise. Quality adjudication should be blinded to the arm, while authorization, provenance and timing evidence determine the other components of the conjunctive endpoint (C16).

Registration before routing should distinguish requested tasks, authorized tasks, tasks permitted to undergo both evaluations, and actual admissions in each arm. The primary rates use the same fixed paired denominator, not only admitted or completed tasks: model mismatch, queue rejection, paid-cap exhaustion, outage, non-delivery and missing required evidence remain outcomes rather than post hoc exclusions. Tasks prohibited from paired evaluation must be reported separately with reasons and arm-specific coverage, never exported merely to make the comparison complete.

Reporting the paired fraction of requested demand makes the efficacy estimand\textquoteright{}s eligibility restriction visible, while the four-baseline analysis retains demand that the primary contrast cannot ethically or contractually include (C16, C19).

Pairing also requires controls against interference and leakage. The controlled design should match load conditions and isolate caches and resource envelopes sufficiently to support its declared sampling model; historical gateway observations cannot substitute for contemporaneous pairs. Independence concerns both task sampling and execution randomness, not merely different prompts or researchers. Shared outages, queues, caches, repeated templates or a diminishing common budget can correlate episodes, in which case live-service observations require a separately prespecified independent-block analysis or descriptive reporting rather than being counted as independent pairs (C16, C18).

At the fixed final analysis, let $X$ count federation-only successes and $Z$ gateway-only successes among $n\geq1$ pairs. Define $p_{10}=P(Y_F=1,Y_G=0)$, $p_{01}=P(Y_F=0,Y_G=1)$ and the useful-completion difference $\Delta=p_{10}-p_{01}$. The proposed conservative procedure forms equal-tailed Clopper\textendash{}Pearson intervals for the two discordance probabilities, each at confidence $0.975$, and uses the relevant endpoints to construct \$\textbackslash{}ell=\textbackslash{}max\textbackslash{}\{-1,l(X)-u(Z)\textbackslash{}\}\$. Marginal binomial coverage and a union bound provide at least $0.95$ coverage for the resulting difference interval under the stated independent, identically distributed episode model, without requiring independence between $X$ and $Z$ (C16).

The efficacy gate is \$\textbackslash{}ell\textgreater{}0.10\$ in one approved, prespecified primary population. A positive point estimate, or an interval that merely excludes zero, does not establish an improvement exceeding ten percentage points; an interval crossing that margin is inconclusive, not evidence of equivalence. Additional confirmatory strata require a prespecified multiplicity adjustment rather than replacing the primary population with a favorable subgroup. Efficacy remains only one gate: policy, accounting, service, fairness and administrative-cost requirements must pass separately, and subgroup coverage must remain visible even when the primary result is favorable (C16, C19).

Neither thirty tasks per stratum nor the earlier approximate $486$-pair calculation justifies enrollment for this exact conservative decision procedure. The saved analytical certificate instead claims that $N=1800$ independent, identically distributed paired episodes suffice for at least $80\%$ probability of passing the efficacy gate under the explicitly hypothetical alternative $\Delta=0.20$ and discordance probability $q=0.40$. Its argument bounds the conservative binomial endpoints using Hoeffding\textquoteright{}s inequality and then bounds failure of a sufficient event using Cantelli\textquoteright{}s inequality.

This remains an unverified analytical planning result, not a measured workload property, minimum sample size, affordability finding or probability of passing all deployment gates; it entails $3600$ arm executions before child calls and separate reproducibility repeats (C17, C18).

The statistical lead should therefore approve the primary population, executable endpoint, independence or block justification, multiplicity plan, missing-outcome treatment and fixed enrollment only after finance and computing certify a feasible design. A smaller affordable sample is not automatically inadequate, because a sharper analysis could justify it, but the earlier approximation cannot certify its power for the actual gate. If statistical adequacy and resource feasibility cannot both be established, the 90-day activity should remain a preparation or descriptive feasibility study rather than be presented as a confirmatory efficacy pilot.

Safety stop/go authority must be assigned in advance, and a stopped or incomplete study must not inherit the completed fixed-sample certificate (C12, C17, C18, C19).

The final report should separate controlled paired inference from live-service generality and retain exclusions, failed admissions, safe-denial availability losses, dependence limitations and administrative effort alongside the completion result. Across all four baselines, it should connect coverage, quality, provenance and version reproducibility, latency, throughput, availability, fairness and costs to the same declared workload and accounting horizon. The evidence thus changes the burden of argument: federation must demonstrate an approved capability or completion gain that justifies its additional costs and trust boundaries, rather than merely show that remote capacity exists.

Unless that case passes the separate gates, the campus gateway with separately approved frontier access remains the provisional smaller-boundary alternative (C10, C12, C16, C19).

\begin{longtable}{@{}>{\raggedright\arraybackslash}p{\dimexpr(\linewidth-2\tabcolsep*2)/126*42\relax}>{\raggedright\arraybackslash}p{\dimexpr(\linewidth-2\tabcolsep*2)/126*42\relax}>{\raggedright\arraybackslash}p{\dimexpr(\linewidth-2\tabcolsep*2)/126*42\relax}@{}}
\caption{Dependency gates for the proposed 90-day evaluation. Each row authorizes progression only after its conditions are met; none records a completed approval or test.}\\
\toprule
\textbf{Gate} & \textbf{Required decision or evidence} & \textbf{Consequence if unresolved}\\
\midrule
\endfirsthead
\toprule
\textbf{Gate} & \textbf{Required decision or evidence} & \textbf{Consequence if unresolved}\\
\midrule
\endhead
\midrule
\multicolumn{3}{r}{\footnotesize Continued on next page}\\
\endfoot
\bottomrule
\endlastfoot
Source and institutional approval & Inspect both starting references with actual access dates and versions; obtain campus approval for processing domains, derived artifacts, retention, revocation, contracts and liability ownership. & Do not treat proposed extensions as verified gateway capabilities or enroll real-data workloads.\\
Executable assessment and fault-test preparation & Approve task rubrics, provenance and deadline rules, fairness reference allocation and administrative ceilings; implement and test mediation, cancellation, recovery and audit behavior. & Retain public or synthetic preparation; do not equate expected fault-test behavior with observed enforcement.\\
Resource-feasibility certification & Finance certifies dated prices, enforceable charge bounds and opening liabilities; computing certifies compatible capacity, a feasible schedule and staff/grading availability. & Do not select enrollment from statistical planning alone or infer deadline feasibility from aggregate GPU-hours.\\
Paired enrollment authorization & Freeze one primary population, comparable entitlements, parent-task endpoint, independence or block model, multiplicity rule, missing-outcome treatment and justified fixed enrollment. & Conduct descriptive feasibility work rather than claim a confirmatory efficacy experiment.\\
Fixed final analysis and decision & Apply the prespecified paired confidence-bound gate and separately assess policy, accounting, service, fairness and administrative-cost gates; preserve four-baseline coverage and cost reporting. & An interval crossing the efficacy margin is inconclusive; a favorable efficacy result cannot override another failed or unassessed gate.\\
\end{longtable}
\par\smallskip{\footnotesize Evidence: [Research session 1, artifact 6], [Research session 1, artifact 7], [Research session 1]; claims: C10, C12, C16, C17, C18, C19}
\textit{Claims discussed: }C10, C12, C16, C17, C18, C19
\section{Algorithms, results, and limitations}
\subsection{C1 -- Limitation}
The supplied characterization attributes client-protocol translation, Ollama/vLLM backends, WDRR scheduling, per-user quotas, audit logging and GPU telemetry to MindRouter's campus gateway. The starting documentation and PEARC'26 paper were not retrieved in the saved analysis, so their current contents and changed claims remain unverified.
The evidence explicitly distinguishes project-supplied descriptions of published capabilities from independently measured facts. The reported review date, 2026-10-31, is historical information supplied with the project, not an access date established by this analysis. Federation and frontier routing cannot be attributed to MindRouter on this basis.
\textit{Evidence: }[Research session 1, artifact 2], [Research session 1]
\subsection{C2 -- Established}
The comparison defines $\mathsf{L}$ as campus-local inference only, $\mathsf{A}$ as approved frontier API only, $\mathsf{G}$ as a campus gateway with separately approved external access, and $\mathsf{F}$ as a hypothetical federation added to that gateway baseline.
These are analytical definitions, not observations of four deployed systems. Separately approved external access does not imply integrated frontier routing, and adding a federation does not relax the home campus's eligibility rules.
\textit{Evidence: }[Research session 1, artifact 5], [Project overview]
\subsection{C3 -- Established}
A useful completion is one parent task that satisfies authorized processing, its prespecified quality rubric, required provenance and its delivery deadline. Offered but unserved tasks remain in coverage denominators, while child calls and retries are resource consumption rather than additional successes.
This operational definition prevents a multi-call workflow from appearing more productive merely because it makes more calls. Correct policy refusal is reported as a safety outcome, not a delivered research result. Model-quality comparisons additionally need matched policy-eligible strata.
\textit{Evidence: }[Research session 1, artifact 2], [Research session 1, artifact 8], [Research session 1]
\subsection{C4 -- Conjecture}
The proposed routing rule constructs a home-campus-authorized destination set before ranking quality, deadline feasibility or incremental cost. Ambiguity produces no eligible destination, and every tool call or destination change requires renewed eligibility checking.
This is an architecture candidate whose enforcement requires mediation of network egress, tools, logging, storage and derived artifacts. Approval must cover identity, classification, destination, retention, contract, audit and quota requirements. A routing recommendation alone cannot establish containment or recipient compliance.
\textit{Evidence: }[Research session 1, artifact 3], [Project overview]
\subsection{C5 -- Conjecture}
Under the stated partition model, a disconnected admission point cannot both admit new requests from cached approval and guarantee immediate enforcement of an unseen home-campus revocation.
The proof candidate compares unchanged-approval and revoked-approval executions with identical requests, cached policy, clocks and local observations. The admission point therefore cannot choose different outcomes based on the unseen revocation; positive admission probability in one execution also permits unauthorized admission in the other. The argument is conditional and not formally verified, applies to new authorization decisions rather than data recall, and does not establish a defect in MindRouter.
\textit{Evidence: }[Research session 1, artifact 1]
\subsection{C6 -- Conjecture}
The proposed admission machine conditionally preserves current-policy authorization at the authoritative dispatch-grant commit and the scoped accounting invariant $L_s+H_s\leq B_s$, where $L_s$ is settled liability, $H_s$ is reserved unresolved liability and $B_s$ is the cap for scope $s$.
The invariant argument assumes complete mediation, trustworthy classification bindings, a linearizable home authority, atomic durable multi-scope accounting, enforceable charge bounds and authoritative settlement. Preliminary approval is advisory; final authorization and grant consumption share the authoritative commit point. Decision-time authorization does not prevent a pre-revocation grant from transmitting later.
\textit{Evidence: }[Research session 1, artifact 4], [Research session 1, artifact 3]
\subsection{C7 -- Established}
Timeout or cancellation alone cannot justify releasing an unresolved reservation: two liabilities of $6$ can exceed a cap of $10$ if the first reservation is released before its liability is resolved.
A timeout does not establish that the first call was uncharged. Releasing its hold permits a second admission even though both charges remain possible, producing exposure of $12$. This hand-worked counterexample refutes premature release, not the conditional reservation invariant.
\textit{Evidence: }[Research session 1, artifact 4], [Research session 1, artifact 7]
\subsection{C8 -- Conjecture}
A proposed burn-before-send adapter consumes a durable dispatch grant before transmission and binds its acknowledgement to one original live continuation. Under the stated sender assumptions, this supports at-most-one adapter invocation initiation per attempt, not exactly-once provider execution.
A one-shot transport latch and refusal to reconstruct sending authority from restarts, status queries, queues or takeover sacrifice availability to avoid replay. These mechanisms remain unimplemented assumptions. Financial coverage must persist through ambiguous outcomes, including a still-possible delayed send.
\textit{Evidence: }[Research session 1, artifact 7]
\subsection{C9 -- Established}
A fresh issuer namespace does not make stale-backup financial recovery safe. If an old liability of $6$ disappears from the restored ledger, a new liability of $6$ still exceeds a cap of $10$, even when replay is prevented.
Namespace fencing controls old sending authority, not missing financial exposure. Recovery must also preserve, reconcile or conservatively reserve prior obligations before new admissions. The counterexample invalidates a recovery shortcut without disproving the conditional machine under its liability-preservation assumption.
\textit{Evidence: }[Research session 1, artifact 7]
\subsection{C10 -- Limitation}
WDRR weights alone do not establish unconditional latency or starvation freedom. Deadline-feasible service also depends on compatible model capacity, blocking, queues, cold starts, maintenance and offered demand.
The saved evidence explicitly retains counterexamples to unconditional scheduling guarantees and identifies model fit, blocking, queues and fairness as measurement obligations. It supplies no executed scheduling benchmark, end-to-end service guarantee or routing scalability result.
\textit{Evidence: }[Research session 1, artifact 3], [Research session 1], [Project overview]
\subsection{C11 -- Established}
For a fixed workload, horizon and cost basis, with $C_G>0$, $S_G>0$, $\Delta C>0$ and $S_G+\Delta S>0$, federation lowers cost per useful completion exactly when $\Delta S/S_G>\Delta C/C_G$.
Comparing $(C_G+\Delta C)/(S_G+\Delta S)$ with $C_G/S_G$ permits cross-multiplication because both completion denominators are positive. Cancelling the shared product gives $S_G\Delta C<C_G\Delta S$, yielding the stated proportional threshold. This accounting identity establishes neither measured savings nor permission to trade away safety or mandatory coverage.
\textit{Evidence: }[Research session 1, artifact 2]
\subsection{C12 -- Unresolved}
A sourced four-baseline total-cost comparison remains incomplete. Current API rates, campus cost inputs, opening liabilities, compatible capacity and active-researcher denominators have not been established sufficiently to report campus-specific monthly or annual totals, break-even points or useful-task costs.
The evidence reports a hand-auditable hypothetical worksheet but does not expose its complete numerical ledger in the bounded index. The required comparison separates sunk, incremental and fully loaded cost and includes local ownership, operations and API workflow amplification. Reservation scenarios are not dated token prices or observed spending.
\textit{Evidence: }[Research session 1], [Research session 1, artifact 5], [Project overview]
\subsection{C13 -- Established}
Aggregate local usefulness of $800/1000$ does not imply useful completion of $160/200$ sensitive tasks. Zero sensitive successes and $800$ other successes satisfy the aggregate figure while failing the sensitive requirement.
The counterexample preserves the total number of local successes while changing their distribution across mandatory strata. It refutes the earlier aggregate-to-sensitive feasibility inference. Giving sensitive tasks admission priority does not itself establish their quality or on-time completion.
\textit{Evidence: }[Research session 1, artifact 5], [Research session 1, artifact 8]
\subsection{C14 -- Established}
In the repaired hypothetical monthly scenario, $200$ sensitive and $800$ other local admissions each have useful-completion probability $0.80$, selected API admissions have probability $0.90$, and $n=\min(1000,\lfloor B/r\rfloor)$ for enforceable per-root reservation $r>0$. Expected useful completions are $800+0.9n$, including $160$ sensitive completions in expectation; the proposed $70\%$ overall expected-count target requires $\lfloor B/r\rfloor\geq667$.
Linearity of expectation gives $200(0.80)+800(0.80)=800$ local successes without requiring independence. The overall target is $1400$ of $2000$ offered roots, so $800+0.9n\geq1400$ requires integer admission count $n\geq667$. These imposed probabilities repair the expectation calculation but provide no realized sensitive-service guarantee.
\textit{Evidence: }[Research session 1, artifact 8], [Research session 1, artifact 5]
\subsection{C15 -- Established}
If federation provides zero usable additional peer capacity, preserves identical service opportunities, yields no offsetting benefit and adds strictly positive overhead on the chosen cost basis, the gateway baseline has the same useful completions at lower cost.
This scoped dominance result follows from equal completion counts and a positive cost increment. Policy exclusions, incompatible capacity or missed deadlines can make nominal peer GPUs unusable, but their actual incidence is unknown. The result supports a counterexample to unconditional federation advantage, not a general recommendation against cooperation.
\textit{Evidence: }[Research session 1, artifact 5], [Research session 1, artifact 8]
\subsection{C16 -- Conjecture}
The proposed primary pilot gate requires a prespecified paired lower confidence bound \$\textbackslash{}ell\textgreater{}0.10\$ for the federation-minus-gateway useful-completion difference in one approved population, together with separate policy, accounting, service, fairness and administrative-cost gates.
The proposal uses paired parent-task episodes, joint success criteria, blinded quality adjudication and a conservative discordance-based interval at a fixed final analysis. It distinguishes evidence of any improvement from evidence exceeding ten percentage points. An interval crossing the margin is inconclusive, not evidence of equivalence.
\textit{Evidence: }[Research session 1, artifact 6], [Research session 1]
\subsection{C17 -- Limitation}
Thirty tasks per stratum and the earlier approximate $486$-pair planning calculation do not establish power for the proposed exact conservative confidence-bound gate.
The saved critique distinguishes approximate normal planning from the operating characteristics of the actual decision procedure. Small-stratum precision was explicitly inadequate in the hypothetical examples. The later analytical certificate addresses a narrower power question but does not make the pilot enrollment-ready.
\textit{Evidence: }[Research session 1, artifact 6], [Research session 1]
\subsection{C18 -- Conjecture}
The saved analytical certificate claims that $N=1800$ independent, identically distributed paired episodes suffice for at least $80\%$ probability of passing the efficacy gate under the hypothetical alternative $\Delta=0.20$ and discordance probability $q=0.40$. This is a sufficient planning bound, not a minimum sample size, affordability result or deployment-success probability.
The indexed account describes bounding conservative binomial endpoints with Hoeffding's inequality and a sufficient-event failure probability with Cantelli's inequality. The bounded preview does not expose the full numerical derivation, so the plan preserves the certificate as an unverified analytical result. Independence concerns parent-task episodes and their execution randomness, not merely distinct prompts or researchers.
\textit{Evidence: }[Research session 1]
\subsection{C19 -- Unresolved}
The broader feasibility question remains unresolved: no source-verified capability review, campus-specific cost and workload evidence, implemented enforcement, approved enrollment protocol or measured federation advantage is supplied.
Institutional acceptance of classification inheritance, decision-time revocation, recipient obligations, contractual liability ceilings and final billing remains an input to feasibility. Conditional proofs and hypothetical worksheets cannot replace those decisions or empirical service measurements.
\textit{Evidence: }[Research session 1], [Research session 1, artifact 8], [Research session 1, artifact 4]
\section{Limitations}
\begin{itemize}
\item The documentation at https://mindrouter.ai/documentation.html and the PEARC'26 paper at https://ai4ra.uidaho.edu/wp-content/uploads/2026/04/Sheneman\_MindRouter\_PEARC26\_Short\_Paper.pdf were not retrieved in the saved analysis. Their reported historical review date does not establish a new access date or verify current capabilities.
\item The bounded index contains truncated extracts and pointer-only relationships, not the complete underlying proofs, numerical ledger or implementation artifacts. Its selection cannot establish exhaustive evidence coverage.
\item Authorization, accounting and one-use dispatch results are conditional proof candidates. Complete mediation, charge ceilings, recipient obligations, recovery correctness and operational release predicates have not been implemented or independently verified.
\item No campus workload traces, quality results, scheduling benchmarks, current API prices, measured administrative burden or empirical four-baseline comparison are available.
\item The repaired worksheet establishes hypothetical expected counts, not realized service guarantees. The pilot certificate concerns a specified efficacy gate under hypothetical sampling assumptions, not affordability or the probability of passing all deployment gates.
\end{itemize}
\section{Open obligations}
\begin{itemize}
\item Retrieve and version both starting sources, record actual access dates, and verify protocol translation, backend support, WDRR, quotas, audit logging, telemetry and campus scope without attributing proposed extensions to MindRouter.
\item Inspect the full original accounting, authorization, dispatch and pilot artifacts, including the pointer-only target, before asserting exact refinement relationships or reproducing omitted numerical derivations.
\item Obtain campus decisions on classification inheritance, student and patient data, controlled and sponsor-restricted material, unpublished research, approved domains, retention, audit, contracts, liability ownership and revocation timing.
\item Implement and fault-test complete mediation, immutable protocol normalization, atomic nested reservations, authoritative grant ordering, one-use issuance, cancellation, sender fencing and crash recovery. Establish a release predicate that excludes remaining delayed-send liability.
\item Populate the local and API cost ledger with dated prices, enforceable root-task liability bounds, opening obligations, compatible capacity, staff costs and active-researcher denominators. Complete monthly and annual, incremental and fully loaded, break-even and sensitivity comparisons.
\item Approve a primary population, executable task-specific rubrics, provenance and deadline rules, fairness reference allocation, sampling or independent-block structure, multiplicity plan, fixed enrollment and administrative ceiling.
\item Check whether the proposed study fits the 90-day execution, grading, capacity and budget envelopes. If the justified enrollment is unaffordable or independence is indefensible, declare the efficacy study infeasible or underpowered rather than relaxing the gate after observing outcomes.
\item Assign campus computing, security, procurement, faculty and student representatives responsibility for service targets, destination approval, billing and retention terms, scientific usefulness, accessibility and appeal procedures. Resolve their questions before enrollment.
\item Run the approved pilot and report safety events, useful-task outcomes, uncertainty, service metrics and administrative costs before claiming reliable access or a federation deployment advantage.
\end{itemize}
\section{Frozen obligation reconciliation}
Unlisted obligations are not presumed resolved. The archived research snapshot retains complete source references and the full reconciliation.
\begin{itemize}
\item Obtain and inspect the original pointer-only target [Research session 1, artifact 3] and full peer artifacts before certifying their exact relationship; retain the supplied extracts' truncation uncertainty. \textit{ (Not listed by author; not presumed resolved; 1 saved source reference(s))}
\item Verify both MindRouter starting references with actual access dates and versions. Published capability claims and changed claims remain pending. \textit{ (Not listed by author; not presumed resolved; 1 saved source reference(s))}
\item Approve the primary population, executable rubrics, provenance and latency rules, multiplicity plan, defensible sampling units, fairness reference allocation and fixed enrollment. The analytical power component is advanced, not a fully resolved enrollment obligation. \textit{ (Not listed by author; not presumed resolved; 1 saved source reference(s))}
\item Populate and sign the symbolic resource constraints with dated prices, enforceable charge bounds, opening liabilities, compatible capacity, independent execution opportunities and staff/grading availability. Exact or sharper power calculations remain optional routes to certifying a smaller affordable design. \textit{ (Not listed by author; not presumed resolved; 1 saved source reference(s))}
\item Obtain campus security/legal/procurement decisions on classification and derived artifacts, approved domains, retention, revocation timing, contracts and liability ownership. Implement and fault-test mediation, grant ordering, one-use issuance, cancellation, ledger recovery and audit controls. \textit{ (Not listed by author; not presumed resolved; 1 saved source reference(s))}
\item Measure model fit, blocking, cold starts, queues, availability and fairness; do not infer starvation freedom from WDRR weights. Preserve the recorded constraints against premature timeout release and unsafe fallback. \textit{ (Not listed by author; not presumed resolved; 1 saved source reference(s))}
\item Complete sourced four-baseline monthly/annual, incremental/fully loaded, break-even and sensitivity analyses with active-researcher and useful-parent-task denominators; run the approved pilot before claiming useful access or federation benefit. \textit{ (Not listed by author; not presumed resolved; 1 saved source reference(s))}
\item Retrieve both starting sources, record actual access dates and versions, and verify the user-supplied characterization of protocol translation, WDRR, quotas, audit logging, telemetry, and on-campus scope. No federation or frontier-routing feature has been verified. \textit{ (Not listed by author; not presumed resolved; 7 saved source reference(s))}
\item Ask campus security and data owners to define classification inheritance, approved domains, revocation timing, in-flight handling, and acceptable lease semantics; ask procurement/legal to validate partner/provider retention, subprocessors, audit rights, and enforceable charge bounds. \textit{ (Not listed by author; not presumed resolved; 7 saved source reference(s))}
\item Implement and independently inspect complete egress mediation, protocol normalization, credential and replay controls, protected logging, and atomic multi-level reservations. Recipient compliance and malicious administrators are outside the proved authorization lemma. \textit{ (Not listed by author; not presumed resolved; 7 saved source reference(s))}
\item Execute the proposed paired partition and adversarial tests. Determine whether synthetic tests can safely progress to a 90-day pilot; establish stop/go authority and measure availability lost to safe denial. \textit{ (Not listed by author; not presumed resolved; 7 saved source reference(s))}
\item Measure task quality, reproducibility, latency, GPU/model fit, utilization, fairness, administrative effort, and fully loaded costs across all four baselines. No numerical break-even, budget sufficiency, or useful-task completion conclusion is supported yet. \textit{ (Not listed by author; not presumed resolved; 7 saved source reference(s))}
\item Test whether federation supplies a material, policy-permitted capability unavailable through campus-local service plus separately approved frontier access; otherwise prefer the smaller trust boundary provisionally. Literature novelty and empirical generality have not been assessed. \textit{ (Not listed by author; not presumed resolved; 7 saved source reference(s))}
\item Retrieve the MindRouter documentation and PEARC'26 paper, record actual access dates and versions, and verify the supplied feature descriptions and any changes. Obtain dated API rates, contractual charging bounds, retention terms and institutional approval; none is established here. \textit{ (Not listed by author; not presumed resolved; 7 saved source reference(s))}
\item Collect a campus cost ledger that separates sunk, incremental and fully loaded cost, including replacement timing, power/cooling, storage/network, support, procurement and staff effort. Replace illustrative aggregate totals before making a procurement recommendation. \textit{ (Not listed by author; not presumed resolved; 7 saved source reference(s))}
\item Measure matched task quality, complete workflow token/call amplification, latency distributions, fit/cold-start behavior, capacity, unserved demand and provenance limits. Estimate uncertainty and sensitivity to task mix, utilization, API prices and demand rather than extrapolating constant success fractions through saturation. \textit{ (Not listed by author; not presumed resolved; 7 saved source reference(s))}
\item Establish whether all paid calls, tools and retries can be assigned enforceable maximum charges and whether timeout reconciliation and nested quotas preserve the reservation invariant under concurrency and failures. This requires implementation testing and contract review. \textit{ (Not listed by author; not presumed resolved; 7 saved source reference(s))}
\item Resolve campus identity, approved-domain classification, derived-artifact handling, retention, audit, incident and cross-campus agreement requirements before any sensitive or federated testing. \textit{ (Not listed by author; not presumed resolved; 7 saved source reference(s))}
\end{itemize}
222 further distinct obligations are retained in the archived reconciliation.
\section{Conclusions}
Federation is not justified by nominal GPU supply or call throughput alone. Under a common workload and cost basis, its proportional useful-completion gain must exceed its proportional cost increase, while independent policy and service requirements still hold.

Home-campus authority and hard spending caps require explicit authorization, reservation, dispatch and reconciliation semantics. The proposed mechanisms expose important failure cases but do not constitute a verified implementation.

A campus gateway with separately approved external access is the necessary practical comparator. It dominates the hypothetical federation in the scoped case of identical service opportunities, no offsetting benefit and positive federation overhead; broader superiority remains unmeasured.

The next defensible step is a source-verified, institutionally approved and resource-feasible pilot with separate sensitive-task outcomes and an explicit inconclusive result. The present evidence supports that evaluation plan, not a deployment recommendation.

\textit{Claims discussed: }C3, C4, C6, C7, C8, C9, C11, C13, C15, C16, C18, C19
\appendix
\section{Selected source-linked research excerpts}
The following are bounded, selected excerpts of original research records. They are not exhaustive; consult the full frozen snapshot archived with this paper. They do not constitute independent verification.
\subsection{Research excerpt 1}
\textit{Source and field: }[Research session 1, artifact 3]: skeptic\_output/reasoning\_assessments\par
[\{"step\_number": 1, "disposition": "unresolved", "rationale": "The attribution matches the supplied project description, and the separation of published gateway claims from proposed extensions is appropriate. No original source establishes those capabilities here.", "objection": "The entry is labeled an assumption, but the published-capability premise is not one of the five explicit conditions of the ledger claim. Retain it as unverified background, not an established premise.

The ledger proof does not require it.", "basis\_artifact\_ids": []\}, \{"step\_number": 2, "disposition": "accepted", "rationale": "As a proposed admission rule, constructing eligibility before ranking prevents a better score or failure from authorizing a prohibited destination. Per-call re-evaluation and mediation of derived artifacts address the stated policy scope.", "objection": "", "basis\_artifact\_ids": []\}, \{"step\_number": 3, "disposition": "accepted", "rationale": "Atomic admission against every applicable scope preserves each inequality.

Attributing the same liability to nested scopes does not imply multiple billing, and separating GPU quotas from money avoids conflating resources.", "objection": "", "basis\_artifact\_ids": []\}, \{"step\_number": 4, "disposition": "accepted", "rationale": "The settlement change is z-r \textless{}= 0; unresolved events leave totals unchanged. Together with the initial invariant and the listed durability and accounting-window assumptions, this proves allocation safety. The explicit cash-flow limitation prevents overclaiming invoice-month control.", "objection": "",\ldots{}

\subsection{Research excerpt 2}
\textit{Source and field: }[Research session 1]: steps/4/output/reasoning\_assessments\par
[\{"step\_number": 1, "disposition": "unresolved", "rationale": "The attribution matches the supplied project description, and the separation of published gateway claims from proposed extensions is appropriate. No original source establishes those capabilities here.", "objection": "The entry is labeled an assumption, but the published-capability premise is not one of the five explicit conditions of the ledger claim. Retain it as unverified background, not an established premise.

The ledger proof does not require it.", "basis\_artifact\_ids": []\}, \{"step\_number": 2, "disposition": "accepted", "rationale": "As a proposed admission rule, constructing eligibility before ranking prevents a better score or failure from authorizing a prohibited destination. Per-call re-evaluation and mediation of derived artifacts address the stated policy scope.", "objection": "", "basis\_artifact\_ids": []\}, \{"step\_number": 3, "disposition": "accepted", "rationale": "Atomic admission against every applicable scope preserves each inequality.

Attributing the same liability to nested scopes does not imply multiple billing, and separating GPU quotas from money avoids conflating resources.", "objection": "", "basis\_artifact\_ids": []\}, \{"step\_number": 4, "disposition": "accepted", "rationale": "The settlement change is z-r \textless{}= 0; unresolved events leave totals unchanged. Together with the initial invariant and the listed durability and accounting-window assumptions, this proves allocation safety. The explicit cash-flow limitation prevents overclaiming invoice-month control.", "objection": "",\ldots{}

\subsection{Research excerpt 3}
\textit{Source and field: }[Research session 1, artifact 1]: skeptic\_output/reasoning\_assessments\par
[\{"step\_number": 1, "disposition": "accepted", "rationale": "Accepted as the exact conditional model, not as established university policy. It explicitly excludes alternative information channels and previously committed request-specific authorization, and defines safety and availability sufficiently for the argument.", "objection": "", "basis\_artifact\_ids": []\}, \{"step\_number": 2, "disposition": "accepted", "rationale": "An unobserved home-campus revocation can be inserted without changing the admission point's request, credentials, clock or cached state.

A signature on historical policy provides authenticity, not knowledge that no later revocation occurred.", "objection": "", "basis\_artifact\_ids": []\}, \{"step\_number": 3, "disposition": "accepted", "rationale": "Identical observations imply identical deterministic decisions. Coupling random choices gives equal admission probability in the two executions, so zero unauthorized-admission probability excludes positive partition availability under these assumptions.", "objection": "", "basis\_artifact\_ids": []\}, \{"step\_number": 4, "disposition": "accepted", "rationale": "Campus-approved non-revocable leases weaken the immediate-revocation requirement, while an authoritative connected decision service removes the disconnected-authorizer premise.

The argument correctly does not establish recall or prevention of delivery after an earlier authorization.", "objection": "", "basis\_artifact\_ids": []\}, \{"step\_number": 5, "disposition": "accepted", "rationale": "Accepted as a proposed conservative policy specification. Tools and\ldots{}

\subsection{Research excerpt 4}
\textit{Source and field: }[Research session 1, artifact 4]: skeptic\_output/reasoning\_assessments\par
[\{"step\_number": 1, "disposition": "accepted", "rationale": "The supplied prior extract identifies the stated policy and reservation gaps. Both peer extracts explicitly list liability-bound and complete-mediation premises. The step correctly treats these as unverified inputs and does not import the truncated economic argument.", "objection": "", "basis\_artifact\_ids": ["[Research session 1, artifact 1]", "[Research session 1, artifact 2]", "[Research session 1, artifact 3]"]\}, \{"step\_number": 2, "disposition": "accepted", "rationale": "Accepted as conditional model definitions.

Separating root task, immutable logical operation and potentially billable attempt prevents retries from becoming extra useful completions and prevents authorization reuse after changing a destination or payload.", "objection": "", "basis\_artifact\_ids": []\}, \{"step\_number": 3, "disposition": "accepted", "rationale": "Accepted as an explicit authority and campus-acceptance premise, not established policy.

Linearization gives revocation and grant a definite order while expressly allowing later transmission from an earlier grant; disconnected cached approval cannot create a new commit.", "objection": "", "basis\_artifact\_ids": []\}, \{"step\_number": 4, "disposition": "accepted", "rationale": "Accepted as the conditional state and trusted-validation specification. Separating financial state from task outcome avoids inferring settlement from success, failure or cancellation. Unknown classifications and unsupported request semantics fail closed rather than being silently dropped.", "objection":\ldots{}

\subsection{Research excerpt 5}
\textit{Source and field: }[Research session 1, artifact 6]: skeptic\_output/reasoning\_assessments\par
[\{"step\_number": 1, "disposition": "accepted", "rationale": "The supplied cycle-1 extract contains the stated reservation and scheduling examples. The current proposal explicitly tightens, rather than retrospectively reinterprets, the earlier federation gate and preserves its open requirements.", "objection": "", "basis\_artifact\_ids": ["[Research session 1, artifact 3]"]\}, \{"step\_number": 2, "disposition": "accepted", "rationale": "As a proposed comparison, matching non-federation entitlements isolates the addition of remote capacity.

Explicit financial mediation and policy-first eligibility avoid treating approval or failure as sufficient authorization.", "objection": "", "basis\_artifact\_ids": ["[Research session 1, artifact 3]"]\}, \{"step\_number": 3, "disposition": "accepted", "rationale": "A common preregistered paired denominator prevents admission and completion filtering from inflating arm success. Reporting requested, authorized, paired and admitted counts exposes eligibility selection and limits the estimand's scope.", "objection": "", "basis\_artifact\_ids": []\}, \{"step\_number": 4, "disposition": "accepted", "rationale": "Success is conjunctive in quality, provenance, authorization and delivery deadline.

End-to-end latency includes workflow delays rather than only generation time. The numerical deadlines and incomplete pass rules are correctly presented as proposals awaiting approval.", "objection": "", "basis\_artifact\_ids": []\}, \{"step\_number": 5, "disposition": "accepted", "rationale": "The pairing, cache isolation and matched-load requirements address\ldots{}

\subsection{Research excerpt 6}
\textit{Source and field: }[Research session 1, artifact 7]: skeptic\_output/reasoning\_assessments\par
[\{"step\_number": 1, "disposition": "unresolved", "rationale": "The evidence boundary and retained decision-time and timeout-release limitations are appropriate. However, this review packet contains pointer-only prior artifacts and ledger summaries rather than the described cycle-2 state-machine extract.", "objection": "The claim that supplied source text provides the earlier states cannot be checked here. Request the relevant definitions only; do not treat a pointer or ledger summary as the original argument.", "basis\_artifact\_ids": []\}, \{"step\_number": 2, "disposition": "accepted", "rationale": "Accepted as a proposed specification premise, not an implemented property.

The records bind identity, destination and frozen content, and the explicit classification premise is necessary because a digest alone cannot establish policy correctness.", "objection": "", "basis\_artifact\_ids": []\}, \{"step\_number": 3, "disposition": "accepted", "rationale": "Accepted as a conditional schema and transaction-domain premise. Unique reservation and settlement keys plus exact accounting units support deterministic admission accounting, provided the stated serialization and complete-scope coverage actually hold.", "objection": "", "basis\_artifact\_ids": []\}, \{"step\_number": 4, "disposition": "accepted", "rationale": "Under the stated atomic transaction premise, commit and pre-send closure cannot both win the same state predicate.

Releasing holds only in the winning closure transaction prevents a partially released committed reservation.", "objection": "", "basis\_artifact\_ids": []\},\ldots{}

\subsection{Research excerpt 7}
\textit{Source and field: }[Research session 1, artifact 2]: skeptic\_output/reasoning\_assessments\par
[\{"step\_number": 1, "disposition": "accepted", "rationale": "The feature attribution matches the supplied project description, and the four baseline definitions preserve the distinction between the published campus gateway and hypothetical extensions. The stated 2026-10-31 date is correctly identified as supplied history, not this session's access date.", "objection": "", "basis\_artifact\_ids": []\}, \{"step\_number": 2, "disposition": "accepted", "rationale": "The binary root-task definition entails S\_j as a count of authorized, timely, rubric-passing completions. Keeping refusals and ineligible demand visible prevents treating policy exclusion as successful service.

Separating quality strata and provenance avoids equating call count or version recording with useful, reproducible work.", "objection": "", "basis\_artifact\_ids": []\}, \{"step\_number": 3, "disposition": "accepted", "rationale": "Accepted as a proposed experimental specification under the listed assumption that pilot deadlines, samples and gates require campus approval. The task-specific rubrics and latency decomposition are suitable candidate controls; no achieved service level or model equivalence is inferred.", "objection": "", "basis\_artifact\_ids": []\}, \{"step\_number": 4, "disposition": "accepted", "rationale": "The annualized ownership formula is dimensionally consistent when power is in kW and electricity is priced per kWh.

Cooling and replacement double-counting are explicitly excluded. Distinguishing consolidated transfers, campus budgets, sunk costs and distinct active-researcher denominators is valid\ldots{}

\subsection{Research excerpt 8}
\textit{Source and field: }[Research session 1, artifact 8]: skeptic\_output/proof\_obligations\par
["Local analytical obligation, supported in cycle 3 steps 2\textendash{}4: show aggregate usefulness cannot establish sensitive coverage and state the additional stratum premise.", "Local analytical obligation, supported in cycle 3 steps 5\textendash{}9: distinguish expectations from realized service and derive the fixed-plan cap threshold.", "Local analytical obligation, supported in cycle 3 steps 10\textendash{}11 and 14: scope dominance, liability accounting and the cost-per-useful-completion threshold.", "Unresolved objective obligation: establish documentation and price provenance, actual institutional permissions and the detailed cost evidence required for a campus-specific comparison.", "Unresolved implementation obligation: demonstrate that contractual bounds, mediation, reservation transitions and authorization ordering satisfy the declared premises.", "Unresolved study-design obligation: complete the requested bounded pilot, operational gates and uncertainty register; later measurements must remain distinct from this proposal."]

\subsection{Research excerpt 9}
\textit{Source and field: }[Research session 1, artifact 6]: skeptic\_output/proof\_obligations\par
["Conditional reservation invariant: supported by cycle 1, artifact [Research session 1, artifact 3], supplied reasoning steps 3\textendash{}4; implementation and contractual instantiation remain open.", "Negative constraints on timeout release and unconditional scheduling guarantees: supported by cycle 1, artifact [Research session 1, artifact 3], supplied reasoning steps 5 and 7.", "Paired estimand and variance: supported analytically by cycle 2, current reasoning step 6.", "Conservative interval coverage and multiplicity control: supported under the stated sampling models by cycle 2, current reasoning steps 7 and 11.", "Illustrative precision limits and independent-block bound: supported by cycle 2, current reasoning steps 8\textendash{}10; these are hand-derived planning results, not executed computations.", "Actual gate-power guarantee, fixed affordable enrollment and defensible sampling units: open in cycle 2.", "Operational approvals, sourced costs, implementation safety and measured useful access: open across cycles 1\textendash{}2."]

\subsection{Research excerpt 10}
\textit{Source and field: }[Research session 1, artifact 7]: skeptic\_output/proof\_obligations\par
["Cycle 3 steps 4\textendash{}7 support the conditional uniqueness argument; implementation must establish their transaction, acknowledgement and one-shot-continuation premises.", "Cycle 3 steps 10, 12\textendash{}14 support conditional coverage during ambiguity; an operational finality predicate must cover delayed sending and all applicable financial scopes.", "Cycle 3 step 9 must be corrected so recovery preserves both replay rejection and unresolved financial exposure.", "Exact refinement correspondence requires the absent cycle-2 definitions; no prior-cycle proof obligation is discharged from its pointer.", "Project achievement additionally requires the source-grounded four-baseline comparison, governance review, transparent dated costs, measurements and bounded pilot retained in the ledger."]

\subsection{Research excerpt 11}
\textit{Source and field: }[Research session 1]: strategy\_runs/2/pending\_investigation/reasoning\_steps\par
[\{"title": "1. Evidence boundary and the planning gap", "argument": "The supplied bounded extract of [Research session 1, artifact 6] specifies the confidence-bound gate, paired estimand and approximate 486-pair planning calculation. Its original was not supplied. Artifact [Research session 1, artifact 3] is only a pointer here: neither its original accounting proof nor its original counterexamples was inspected. The ledger records all relevant source, approval, implementation and empirical obligations as unresolved. This cycle supplies a conservative mathematical planning certificate instead of treating the earlier normal approximation as exact power.

It does not establish an affordable enrollment target.", "epistemic\_status": "deduction", "basis\_artifact\_ids": ["[Research session 1, artifact 6]"]\}, \{"title": "2. Independent sampling unit and success definition", "argument": "The sampling unit is one independently sampled parent-task test episode containing both an F execution and a G execution. Success Y\_b=1 requires the frozen quality rubric, authorized processing, required provenance and delivery deadline; otherwise Y\_b=0. Child calls, retries and reproducibility repeats are not additional units. Assume episodes are identically distributed in the declared population and independent in their task and execution randomness.

A defensible controlled design would independently sample workload episodes and isolate their resource envelopes and execution environments. Merely choosing different prompts or researchers does not establish independence. Shared outages,\ldots{}

\subsection{Research excerpt 12}
\textit{Source and field: }[Research session 1, artifact 2]: skeptic\_output/proof\_obligations\par
["Established analytically, agent-reviewed only: the conditional federation threshold and completion-ceiling corollary \textemdash{} cycle 1, current Investigator step de6fe0d9-9414-403f-8c1f-4832ea3e5e5a, reasoning step 6.", "Established analytically within its stated range: the success-adjusted local/API crossover \textemdash{} cycle 1, the same Investigator step, reasoning step 5.", "Established by hand arithmetic as a synthetic example only: greater useful throughput need not lower unit cost \textemdash{} cycle 1, the same Investigator step, reasoning step 7.", "Established conditionally, not operationally: the reservation accounting invariant \textemdash{} cycle 1, the same Investigator step, reasoning step 8; implementation and contract premises remain o18.", "Overall completion still requires o15\textendash{}o21 and an explicit four-baseline cost/coverage sensitivity worksheet; the proposed architecture also needs concrete enforcement transitions."]

\subsection{Research excerpt 13}
\textit{Source and field: }[Research session 1, artifact 3]: skeptic\_output/proof\_obligations\par
["Conditional allocation invariant: supported by cycle 1, Investigator step 0d553ed5-938d-411e-a747-977e366fcaf0, reasoning steps 3\textendash{}4; convincing but unverified.", "Necessity of retaining unresolved reservations: supported by the hypothetical counterexample in cycle 1, reasoning step 5.", "Limits of fair-share scheduling alone: supported by the constructed blocking, incompatibility and strict-priority examples in cycle 1, reasoning step 7.", "Instantiation of accounting and policy assumptions: open; requires contractual evidence and implementation fault tests.", "Decision-ready pilot and empirical usefulness, reliability, cost and federation benefit: open."]

7 further research records are retained in the archived source snapshot.
\section*{Evidence and authorship note}
Snapshot SHA-256: 552f9458e600d3e7d0b3aa6c1dd105cc62dab4967db652123d2f134783b76f34
Generated by an AI author from the cited saved research snapshot. This document is not independent verification. No external bibliography was added.
\section*{Validation warnings}
\begin{itemize}
\item Evidence indexing omitted 9718 material saved research field fragment(s). Those fragments were not reviewed for this paper; the interpretation is provisional where the omitted evidence may matter. Consult the frozen source snapshot and omission ledger for full details.
\item Author self-review: the claim identifiers referenced in narrative linkage fields use zero-padded forms such as C02, whereas the claim register uses forms such as C2. These links require normalization before downstream drafting.
\item Author self-review: several section-purpose strings in this plan are incomplete and must be completed before the plan is used for drafting.
\item Published MindRouter capability claims remain unsupported by direct inspection of the requested primary sources. No current access date, changed feature claim or independent performance verification can be supplied.
\item The original aggregate-to-sensitive feasibility inference is contradicted by a concrete stratum counterexample. Only the repaired conditional expectation calculation should enter the findings.
\item The fresh-namespace recovery shortcut is contradicted by the stale-ledger liability example. Replay prevention cannot substitute for preserving prior financial exposure.
\item The earlier approximate sample-size calculation is not a power calculation for the conservative final gate. The later sufficient certificate is preserved as unverified because its complete derivation is absent from the bounded preview.
\item No result establishes exactly-once provider execution, immediate post-revocation non-delivery, remote retention compliance, unconditional WDRR latency or starvation freedom, or deployment-scale routing complexity.
\item The USD 4,000 monthly cap and USD 2,000/USD 8,000 sensitivities are requested scenarios, not observed campus expenditure. Reservation limits must not be reported as API prices or actual charges.
\item The evidence index explicitly omitted 9718 saved field fragment(s); see the generation coverage ledger and review the complete archived snapshot.
\item Structured paper indexing selected 372 source-linked text fragments from 13 archived sources without AI extraction calls. The complete frozen snapshot remains archived; this bounded index does not constitute review of every saved field.
\item C2 is labeled established, but cited saved evidence includes unverified, incomplete, or conjectural material.
\item C2 is labeled established while the frozen evidence contains counterexample or contradiction language; inspect the cited sources before accepting it.
\item C3 is labeled established, but cited saved evidence includes unverified, incomplete, or conjectural material.
\item C3 is labeled established while the frozen evidence contains counterexample or contradiction language; inspect the cited sources before accepting it.
\item C7 is labeled established, but cited saved evidence includes unverified, incomplete, or conjectural material.
\item C7 is labeled established while the frozen evidence contains counterexample or contradiction language; inspect the cited sources before accepting it.
\item C9 is labeled established, but cited saved evidence includes unverified, incomplete, or conjectural material.
\item C9 is labeled established while the frozen evidence contains counterexample or contradiction language; inspect the cited sources before accepting it.
\item C11 is labeled established, but cited saved evidence includes unverified, incomplete, or conjectural material.
\item A diagram's graph could not be laid out without likely crossings or overlap; its complete nodes, boundaries, and relationships are shown in a readable text representation.
\end{itemize}
\section*{Typesetting verification note}
Needs review: TeX verification is incomplete. The authored research content and safe accepted typesetting corrections were retained. Review the saved proofreading diagnostics before relying on the layout.
\end{document}
